\documentclass[11pt,a4paper]{article}
\usepackage{fontspec}
\usepackage[margin=2.5cm]{geometry}
\usepackage{amsmath,amssymb,amsthm,mathtools,bm}
\usepackage{booktabs}
\usepackage{array}
\usepackage{enumitem}
\usepackage{microtype}
\usepackage[dvipsnames]{xcolor}
\usepackage{tcolorbox}
\usepackage[hidelinks,bookmarksnumbered]{hyperref}
\usepackage{graphicx}

\setlength{\parindent}{0pt}
\emergencystretch=3em
\setlength{\parskip}{0.55em}
\numberwithin{equation}{section}

\newtcolorbox{result}[1][]{colback=Blue!3,colframe=Blue!45!black,boxrule=0.5pt,arc=1.5pt,
  left=6pt,right=6pt,top=4pt,bottom=4pt,title={#1},fonttitle=\bfseries\small,coltitle=white}
\newtcolorbox{remark}{colback=black!3,colframe=black!35,boxrule=0.4pt,arc=1.5pt,left=6pt,right=6pt,top=3pt,bottom=3pt}

\newcommand{\E}{\mathbb{E}}
\newcommand{\Var}{\operatorname{Var}}
\newcommand{\tavg}[1]{\left\langle #1 \right\rangle}
\newcommand{\dd}{\mathrm{d}}
\newcommand{\fn}{f_{\mathrm{n}}}
\newcommand{\wn}{\omega_{\mathrm{n}}}
\newcommand{\wdm}{\omega_{\mathrm{d}}}
\newcommand{\ii}{\mathrm{j}}
\newcommand{\code}[1]{\texttt{#1}}
\newcommand{\Tl}{T_{\mathrm{L}}}

\title{\textbf{vibsynth 2.1 --- Theory manual}\\[0.4em]
\large Mission synthesis, kurtosis control and MMVSD,\\ explained step by step}
\author{Emanuele Pesaresi}
\date{Version 2.1.0, October 2026}

\begin{document}
\maketitle

\begin{abstract}
\noindent This manual explains every concept used by the \code{vibsynth} package, from the definition of a
power spectral density to the Moment-Matched Variable Spectral Density (MMVSD) method introduced in
version~2.1. Each quantity is defined before it is used, and every derivation is written out in full,
so that the text can be followed by reading alone. Part~I builds the common ground (random signals, the
single-degree-of-freedom oscillator, fatigue damage, the Fatigue Damage Spectrum). Part~II describes the
mission synthesis and kurtosis control methods of version~2.0. Part~III derives the MMVSD method, which
prescribes, in addition to the Fatigue Damage Spectrum, the kurtosis of the response of every
oscillator, and its treatment of sine-on-random blocks. The appendices contain the longer
integrations (Lalanne's operator, the ramp-invariant oscillator filter) and the formulas quoted from
the literature.
\end{abstract}

\tableofcontents
\clearpage

\section{Purpose, scope and notation}\label{sec:intro}

\subsection{What the package does}

A component mounted on a vehicle, a machine or a package is shaken during its life. The acceleration
at its base is measured (a \emph{record}), or prescribed by a \emph{specification}, or described by a
\emph{mission}: a list of situations (road types, operating conditions) with the time spent in each.
An accelerated vibration test reproduces, on a shaker and in a shorter time $T$, the fatigue damage
that the life of duration $\Tl$ would cause. \code{vibsynth} computes the acceleration time history
$x(t)$ of such a test (the \emph{test profile}).

The fatigue damage is described by the Fatigue Damage Spectrum (FDS, Section~\ref{sec:fds}): the
damage that the vibration causes in a family of single-degree-of-freedom oscillators with different
natural frequencies. Two signals with the same FDS damage every lightly damped component in the same
way, whatever its natural frequency. Every synthesis method of the package (Parts~II and~III)
therefore starts from the target FDS and builds a signal with that FDS; the methods differ in what
they prescribe besides the FDS: the kurtosis of the signal, its RMS value, its power spectral density,
tonal components, or, in MMVSD, the kurtosis of the response of every oscillator. The \code{edit}
method of Part~IV takes the opposite route: it does not synthesise but keeps the damaging stretches of
the record itself, which makes it right for every S--N exponent, up to a limit that the record sets.

\subsection{How to read this manual}

The manual is written so that it can be followed without paper and pencil. Each symbol is defined at
the place where it first appears, and each derivation shows all intermediate steps. Results that are
used later are collected in boxes like this one:

\begin{result}[Example of a result box]
A boxed equation is a result that later sections refer to.
\end{result}

Three results are taken from the literature without derivation, because they are an empirical fit or
classical theorems: Dirlik's distribution of rainflow ranges, the central limit theorem and the
statistical properties of Welch's PSD estimate (Appendix~\ref{app:quoted}). They are stated precisely
and their role is explained. Everything else is derived here.

\subsection{Notation}

\begin{center}
\begin{tabular}{>{$}l<{$} p{0.72\textwidth}}
\toprule
t & time [s] \\
f,\ \omega = 2\pi f & frequency [Hz] and angular frequency [rad/s] \\
x(t) & base acceleration: the record or the test profile [m/s$^2$] \\
f_s,\ \Delta t = 1/f_s & sampling frequency and sampling interval \\
N,\ T = N\Delta t & number of samples and duration of a signal \\
\Tl & life represented by a target (the duration of use that the test must reproduce) \\
\E[\cdot] & expectation (average over all realisations of a random process) \\
\tavg{\cdot} & time average over the duration of a signal, $\tavg{g} = \frac1T\int_0^T g(t)\,\dd t$ \\
\sigma^2,\ \sigma & variance and standard deviation (RMS value of a zero-mean signal) \\
\kappa & kurtosis (Pearson definition, 3 for a Gaussian variable); $\kappa_z(\fn)$ is called
  \code{k\_z} in the program \\
G(f) & one-sided power spectral density (PSD) [(m/s$^2$)$^2$/Hz] \\
\fn,\ \wn = 2\pi\fn & natural frequency of an oscillator \\
Q,\ \zeta = 1/(2Q) & quality factor and damping ratio \\
z(t) & relative displacement of an oscillator (base minus mass) [m] \\
H(f) & frequency response from base acceleration to relative displacement \\
b,\ C,\ K & Basquin exponent, Basquin constant, stress per unit relative displacement \\
D(\fn) & Fatigue Damage Spectrum (FDS) \\
\mathrm{MRS}(\fn) & Maximum Response Spectrum \\
\Gamma(\cdot) & Euler's Gamma function, $\Gamma(s) = \int_0^\infty u^{s-1}e^{-u}\,\dd u$ \\
\bottomrule
\end{tabular}
\end{center}

The natural frequencies at which the FDS is evaluated form a grid $f_{\mathrm{n},1} <
f_{\mathrm{n},2} < \dots < f_{\mathrm{n},n}$ (the \emph{nodes}). By default the grid has 12 points per
octave: $f_{\mathrm{n},i} = f_{\min}\,2^{(i-1)/12}$, the last node being moved exactly to $f_{\max}$.
A linear grid with a constant step is also available. An index $i$ (or $j$) always runs over the
nodes, so that a quantity defined at every node, such as $D_i = D(f_{\mathrm{n},i})$, is a vector of
length $n$.

\subsection{Organisation}

Part~I (Sections~\ref{sec:signals}--\ref{sec:damage}) introduces random signals, the oscillator and
fatigue damage. Part~II (Sections~\ref{sec:synthesis}--\ref{sec:methods}) explains how a profile with
a prescribed FDS is designed, synthesised and corrected, and the methods of version~2.0. Part~III
(Sections~\ref{sec:mmvsdmodel}--\ref{sec:lines}) derives MMVSD and its treatment of
sine-on-random blocks. Section~\ref{sec:practice}
relates the theory to the settings of the program, and Section~\ref{sec:diagnostics} defines the
diagnostics and checks of the interface. Part~IV (Sections~\ref{sec:limit}--\ref{sec:nsm}, added in
version~2.2 after the practical sections so that the numbering of versions 2.0 and 2.1 is unchanged)
derives the compression limit of a record and the \code{edit} method, the kurtosis a compressed test
should carry, and the non-stationarity matrix behind the response kurtosis spectrum.

\part*{Part I --- Signals, oscillators and damage}
\addcontentsline{toc}{part}{Part I --- Signals, oscillators and damage}

\section{Random signals}\label{sec:signals}

\subsection{Statistics of a signal}

Let $x_1,\dots,x_N$ be the samples of a signal. Its \emph{mean} and its \emph{central moments} of
order $m$ are
\begin{equation}
  \bar x = \frac1N\sum_{k=1}^N x_k, \qquad \mu_m = \frac1N\sum_{k=1}^N (x_k-\bar x)^m .
\end{equation}
The \emph{variance} is $\sigma^2 = \mu_2$, the standard deviation is $\sigma=\sqrt{\mu_2}$, and for a
zero-mean signal $\sigma$ is the root mean square (RMS) value. The \emph{skewness} $\mu_3/\sigma^3$ is
zero for a signal whose values are distributed symmetrically about the mean. The \emph{kurtosis}
(Pearson's definition, used everywhere in this manual and in the program) is
\begin{equation}\label{eq:kurtosis}
  \kappa = \frac{\mu_4}{\mu_2^{2}} .
\end{equation}
It measures the weight of the tails of the distribution of the values: rare large values raise
$\mu_4$ much more than $\mu_2$. The \emph{crest factor} is $\max_k|x_k-\bar x|/\sigma$.

The same definitions apply to a random variable $X$ when the averages over the samples are replaced
by expectations: $\mu_m = \E[(X-\E X)^m]$.

\subsection{The Gaussian distribution and its kurtosis}\label{sec:gausskurt}

A zero-mean Gaussian random variable $X$ with variance $\sigma^2$ has the probability density
\begin{equation}
  p(x) = \frac{1}{\sqrt{2\pi}\,\sigma}\, e^{-x^2/(2\sigma^2)} .
\end{equation}
Its even moments follow from one integration by parts. For an integer $m\ge 0$, write
$x^{m+2}e^{-x^2/(2\sigma^2)} = x^{m+1}\cdot\bigl(x\,e^{-x^2/(2\sigma^2)}\bigr)$ and note that
$x\,e^{-x^2/(2\sigma^2)}$ is the derivative of $-\sigma^2 e^{-x^2/(2\sigma^2)}$. Then
\begin{align}
  \E[X^{m+2}] &= \frac{1}{\sqrt{2\pi}\sigma}\int_{-\infty}^{\infty} x^{m+1}\, x\, e^{-x^2/(2\sigma^2)}\,\dd x
  \nonumber\\
  &= \frac{1}{\sqrt{2\pi}\sigma}\Bigl[-\sigma^2 x^{m+1} e^{-x^2/(2\sigma^2)}\Bigr]_{-\infty}^{\infty}
   + \frac{1}{\sqrt{2\pi}\sigma}\int_{-\infty}^{\infty} (m+1)x^{m}\,\sigma^2 e^{-x^2/(2\sigma^2)}\,\dd x
  \nonumber\\
  &= 0 + (m+1)\,\sigma^2\,\E[X^m] .
\end{align}
The bracket vanishes because the exponential decays faster than any power grows. Starting from
$\E[X^0]=1$: $\E[X^2] = 1\cdot\sigma^2\cdot 1 = \sigma^2$ and $\E[X^4] = 3\sigma^2\E[X^2] = 3\sigma^4$.
Therefore
\begin{result}[Kurtosis of a Gaussian variable]
\begin{equation}\label{eq:gauss4}
  \E[X^4] = 3\sigma^4, \qquad \kappa = \frac{3\sigma^4}{(\sigma^2)^2} = 3 .
\end{equation}
\end{result}
A signal with $\kappa>3$ is called \emph{leptokurtic}: it has more large values than a Gaussian
signal with the same RMS. Measured road and transport vibration is typically leptokurtic, with
$\kappa$ between 4 and 8, because its level changes with the road surface and contains shocks.

\subsection{Power spectral density}\label{sec:psd}

A record of duration $T$ sampled at $f_s$ can be written exactly as a finite Fourier series. With
$\Delta f = 1/T$ and the frequencies $f_m = m\,\Delta f$, $m = 0,1,\dots,N/2$,
\begin{equation}\label{eq:fourier}
  x(t_k) = \bar x + \sum_{m=1}^{N/2} A_m \cos(2\pi f_m t_k + \varphi_m),\qquad t_k = k\Delta t ,
\end{equation}
where $A_m\ge0$ is the amplitude and $\varphi_m$ the phase of the line $f_m$. A cosine of amplitude
$A_m$ has mean square value $A_m^2/2$ over a whole number of periods (because $\cos^2\theta =
\tfrac12 + \tfrac12\cos2\theta$ and the second term averages to zero), and two lines of different
frequencies are orthogonal over the record (the average of their product is zero). The variance of
the record is therefore the sum of the powers of its lines,
\begin{equation}\label{eq:parseval}
  \sigma^2 = \sum_{m} \frac{A_m^2}{2}.
\end{equation}
The \emph{one-sided power spectral density} (PSD) distributes this power over frequency: the power of
the line $f_m$ is spread over the band $\Delta f$ that the line represents,
\begin{equation}\label{eq:psddef}
  G(f_m) = \frac{A_m^2}{2\,\Delta f}, \qquad \sigma^2 = \sum_m G(f_m)\,\Delta f \;\approx\; \int_0^{f_s/2} G(f)\,\dd f .
\end{equation}
The unit of $G$ is (m/s$^2$)$^2$/Hz. For a random process the PSD is the expectation of this quantity
when $T\to\infty$; for a single record it must be \emph{estimated}, because the amplitudes of
individual lines scatter strongly from one record to another. The program uses Welch's estimate:
the record is cut into segments of length $1/r$ ($r$ is the requested resolution in Hz) overlapping
by 50\,\%, each segment is multiplied by a Hamming window (a raised-cosine taper that reduces the
leakage of power between lines), the squared moduli of their Fourier
transforms are scaled as in~\eqref{eq:psddef} and averaged. Averaging many segments reduces the
scatter; this is the only property of the estimate that the manual uses.

\subsection{Gaussian signals with a prescribed PSD}\label{sec:multisine}

Equation~\eqref{eq:fourier} also tells how to \emph{build} a signal with a prescribed PSD $G(f)$:
choose the amplitudes
\begin{equation}\label{eq:amplitudes}
  A_m = \sqrt{2\,G(f_m)\,\Delta f}
\end{equation}
and phases $\varphi_m$ drawn independently and uniformly in $[0,2\pi)$. This is a
\emph{random-phase multisine}. Its variance is $\sum_m A_m^2/2 = \sum_m G(f_m)\Delta f$, exactly the
area under the prescribed PSD, for every choice of the phases. The value $x(t)$ at a given time is a
sum of many independent terms $A_m\cos(2\pi f_m t+\varphi_m)$, each with zero mean and bounded; by the
central limit theorem the distribution of such a sum approaches a Gaussian distribution when many
lines carry comparable power. In practice a test of 60\,s has a line spacing of 1/60\,Hz, and
hundreds of lines fall in any band of interest, so the multisine is Gaussian for all purposes.

A second construction is exactly Gaussian for any number of lines. Write each line as
$A_m\cos(2\pi f_m t+\varphi_m) = U_m\cos(2\pi f_m t) - V_m\sin(2\pi f_m t)$ with $U_m =
A_m\cos\varphi_m$ and $V_m = A_m\sin\varphi_m$ (this is the addition formula of the cosine). If $U_m$
and $V_m$ are drawn as independent zero-mean Gaussian variables with variance $G(f_m)\Delta f$, then
$x(t)$ is, at every $t$, a sum of Gaussian variables with fixed coefficients, hence Gaussian, with
variance $\sum_m G(f_m)\Delta f\,(\cos^2+\sin^2) = \sum_m G(f_m)\Delta f$. The amplitude $A_m =
\sqrt{U_m^2+V_m^2}$ is then random instead of fixed: it is the length of a vector with two independent
Gaussian components, whose distribution (the Rayleigh distribution) is derived in
Section~\ref{sec:rayleigh},
with $\E[A_m^2/2] = G(f_m)\Delta f$. The pair $(U_m, V_m)$, written as the complex number $U_m+\ii V_m$,
is a \emph{complex normal coefficient}; the program draws such coefficients with unit mean square
modulus and scales them to the power of each line.

Keeping the phases (or the coefficients) of a realisation fixed while changing the amplitudes
$A_m$ is called \emph{frozen phases}: the signal then changes smoothly with its PSD, which is what
the closed-loop corrections of Section~\ref{sec:closedloop} rely on.

\subsection{Modulated Gaussian signals}\label{sec:modulated}

A leptokurtic signal can be built by multiplying a stationary Gaussian \emph{carrier} $g(t)$ (zero
mean, variance $\sigma_g^2$) by a positive \emph{envelope} $e(t)$ that varies slowly compared with the
carrier:
\begin{equation}\label{eq:envcarrier}
  x(t) = e(t)\,g(t).
\end{equation}
Divide the duration $T$ into windows short enough for $e$ to be practically constant in each
(window $w$ has duration $T_w$ and envelope value $e_w$) but long enough to contain many periods of
the carrier, so that the time average of $g^2$ and $g^4$ over a window equals its expectation.
Using~\eqref{eq:gauss4} for the carrier,
\begin{align}
  \tavg{x^2} &= \sum_w \frac{T_w}{T}\, e_w^2\, \sigma_g^2 = \sigma_g^2\,\tavg{e^2}, \\
  \tavg{x^4} &= \sum_w \frac{T_w}{T}\, e_w^4\, 3\sigma_g^4 = 3\sigma_g^4\,\tavg{e^4},
\end{align}
and the kurtosis~\eqref{eq:kurtosis} of the product is
\begin{result}[Kurtosis of a modulated Gaussian signal]
\begin{equation}\label{eq:kmod}
  \kappa = \frac{3\sigma_g^4\tavg{e^4}}{\bigl(\sigma_g^2\tavg{e^2}\bigr)^2} = 3\,\frac{\tavg{e^4}}{\tavg{e^2}^2}
  = 3\,\frac{\tavg{s^2}}{\tavg{s}^2}, \qquad s(t) = e(t)^2\sigma_g^2 .
\end{equation}
\end{result}
The quantity $s(t)$ is the \emph{local variance} of the signal. The last form of~\eqref{eq:kmod}
states a general rule used throughout Part~III: if, given a slowly varying local variance $s(t)$,
the signal is Gaussian, its kurtosis is three times the ratio of the mean square of $s$ to the square
of its mean. By the Cauchy--Schwarz inequality, $\tavg{s}^2 \le \tavg{s^2}\tavg{1}$, so
$\kappa\ge3$, with equality only for a constant $s$. Modulation can raise the kurtosis but never lower
it below the Gaussian value.

\section{The single-degree-of-freedom oscillator}\label{sec:oscillator}

\subsection{Equation of motion}

A mass $m$ is connected to a moving base by a spring of stiffness $k_s$ and a viscous damper of
coefficient $c$. The base moves with displacement $y(t)$, so its acceleration is $\ddot y(t) =
x(t)$, the signal of interest. Let $u(t)$ be the absolute displacement of the mass and
\begin{equation}
  z(t) = y(t) - u(t)
\end{equation}
the \emph{relative displacement} (base minus mass). The spring and the damper act on the mass with
the force $k_s(y-u) + c(\dot y-\dot u) = k_s z + c\dot z$, so Newton's law reads $m\ddot u = k_s z +
c\dot z$. Since $u = y - z$, $\ddot u = x - \ddot z$, and
\begin{equation}
  m(x - \ddot z) = k_s z + c \dot z
  \quad\Longrightarrow\quad
  \ddot z + \frac{c}{m}\dot z + \frac{k_s}{m} z = x(t).
\end{equation}
With the \emph{natural angular frequency} $\wn = \sqrt{k_s/m} = 2\pi\fn$ and the \emph{damping
ratio} $\zeta = c/(2m\wn)$,
\begin{result}[Oscillator under base acceleration]
\begin{equation}\label{eq:sdof}
  \ddot z + 2\zeta\wn\dot z + \wn^2 z = x(t), \qquad Q = \frac{1}{2\zeta}.
\end{equation}
\end{result}
The quality factor $Q$ is the usual way to give the damping in vibration testing; $Q=10$ ($\zeta =
5\,\%$) is a common default. The stress in a component modelled by the oscillator is proportional to
the deformation of its spring, $S = K z$ with a constant $K$. The \emph{pseudo-acceleration}
$\wn^2 z$ is often plotted instead of $z$: at low frequency it equals the base acceleration.

\subsection{Frequency response}

Let the base acceleration be the complex exponential $x(t) = e^{\ii\omega t}$. The steady-state
response has the same time dependence, $z(t) = H(\omega)e^{\ii\omega t}$, with $\dot z = \ii\omega z$
and $\ddot z = -\omega^2 z$. Substituting in~\eqref{eq:sdof}:
\begin{equation}
  \bigl(-\omega^2 + 2\ii\zeta\wn\omega + \wn^2\bigr)H(\omega) = 1
  \quad\Longrightarrow\quad
  H(\omega) = \frac{1}{\wn^2-\omega^2 + 2\ii\zeta\wn\omega},
\end{equation}
and, since $|a+\ii b|^2 = a^2+b^2$,
\begin{result}[Frequency response]
\begin{equation}\label{eq:H2}
  |H(\omega)|^2 = \frac{1}{(\wn^2-\omega^2)^2 + (2\zeta\wn\omega)^2}.
\end{equation}
\end{result}
At $\omega=0$, $|H|=1/\wn^2$ (static deflection). At $\omega = \wn$ the first term of the
denominator vanishes and $|H(\wn)| = 1/(2\zeta\wn^2) = Q/\wn^2$: the resonance amplifies the static
response by $Q$.

\emph{Half-power bandwidth.} Near the resonance put $\omega = \wn(1+\varepsilon)$ with
$|\varepsilon|\ll1$. Then $\wn^2-\omega^2 = -\wn^2(2\varepsilon+\varepsilon^2) \approx -2\varepsilon\wn^2$
and $2\zeta\wn\omega\approx 2\zeta\wn^2$, so $|H|^2 \approx 1/\bigl(4\wn^4(\varepsilon^2+\zeta^2)\bigr)$.
This is half of the peak value $1/(4\zeta^2\wn^4)$ when $\varepsilon^2+\zeta^2 = 2\zeta^2$, i.e.\
$\varepsilon = \pm\zeta$. The band between the two half-power points has the width
$2\zeta\wn$, or in hertz
\begin{equation}\label{eq:bandwidth}
  \Delta f = 2\zeta\fn = \frac{\fn}{Q}.
\end{equation}
An oscillator with $Q=10$ responds mainly to a band of $\pm5\,\%$ around $\fn$, i.e.\ about $0.14$
octave wide. With 12 nodes per octave, adjacent oscillators ($2^{1/12}-1\approx6\,\%$ apart) share most
of their band. This overlap appears repeatedly in the design problems of Parts~II and~III.

\subsection{Response to a random input}\label{sec:respvar}

Let $x$ be a stationary random signal with PSD $G(f)$, represented by the lines of
Section~\ref{sec:psd}. Each line $A_m\cos(2\pi f_m t+\varphi_m)$ produces the response line
$|H(2\pi f_m)|A_m\cos(2\pi f_m t+\varphi_m+\arg H)$, so the response has the line powers
$|H|^2A_m^2/2$, and by~\eqref{eq:psddef} its PSD is
\begin{equation}\label{eq:Gz}
  G_z(f) = |H(2\pi f)|^2\, G(f).
\end{equation}
Its variance and, more generally, its \emph{spectral moments} are
\begin{equation}\label{eq:moments}
  m_k = \int_0^\infty f^k\, G_z(f)\,\dd f,\qquad k = 0,1,2,4,\qquad \sigma_z^2 = m_0 .
\end{equation}
Differentiation multiplies a line by $\ii\omega$, so the velocity $\dot z$ has the PSD
$(2\pi f)^2G_z$ and the variance $\sigma_{\dot z}^2 = (2\pi)^2 m_2$; likewise $\sigma_{\ddot z}^2 =
(2\pi)^4 m_4$. These moments describe how often the response crosses zero and how many peaks it has
(Section~\ref{sec:rice}).

\emph{White noise.} If $G(f)=G_0$ is constant, $\sigma_z^2 = G_0\int_0^\infty |H|^2\dd f$. Appendix~A
shows that $\int_0^\infty|H|^2\dd f = 1/(8\zeta\wn^3)$, hence
\begin{equation}\label{eq:miles}
  \sigma_z^2 = \frac{G_0}{8\zeta\wn^3} = \frac{G_0\,Q}{4\wn^3},
  \qquad
  \sigma_{\wn^2 z} = \sqrt{\wn^4\sigma_z^2} = \sqrt{\tfrac{\pi}{2}\,\fn\, Q\, G_0},
\end{equation}
the second form being Miles' equation for the RMS pseudo-acceleration (use $\wn = 2\pi\fn$:
$\wn^4 G_0Q/(4\wn^3) = 2\pi\fn G_0 Q/4 = \tfrac{\pi}{2}\fn QG_0$).

\subsection{The response variance is linear in the PSD: Lalanne's operator}\label{sec:lalanne}

In the program a PSD is described by its values $G_j$ at the nodes $f_{\mathrm{n},j}$ and is linear
between consecutive nodes. Because~\eqref{eq:moments} is linear in $G$, the response variance of the
oscillator $i$ is a linear combination of the nodal values,
\begin{result}[Lalanne's operator]
\begin{equation}\label{eq:W}
  \sigma_{z,i}^2 = \sum_{j=1}^n W_{ij}\,G_j ,
  \qquad W_{ij} = \int_0^\infty |H_i(f)|^2\,\phi_j(f)\,\dd f ,
\end{equation}
\end{result}
where $\phi_j$ is the ``hat'' function equal to 1 at $f_{\mathrm{n},j}$, 0 at the other nodes and
linear in between, and $H_i$ is the response of the oscillator $f_{\mathrm{n},i}$. Lalanne
integrated~\eqref{eq:W} in closed form; Appendix~\ref{app:lalanne} gives the complete derivation. The
matrix $W$ is computed once per damage model. Two details of the program matter:
\begin{itemize}[nosep]
\item the PSD of a synthesised profile is not cut at the first and last node but held constant over a
margin of $3/Q$ relative bandwidth beyond them (an oscillator at the edge of the grid would otherwise
see only half of its band); the methods of version~2.0 leave the effect of this margin to the closed
loop, while the operators of MMVSD include it (Section~\ref{sec:operators});
\item when the pre-filter of Section~\ref{sec:digital} is active, $|H_i|^2$ is multiplied by its power
response.
\end{itemize}

\subsection{How fast the response follows a change of the input}\label{sec:envelope}

Part~III needs to know how the response variance follows an input whose level changes with time.
The mechanical energy of the oscillator per unit mass is $E = \tfrac12\dot z^2 + \tfrac12\wn^2z^2$.
Its derivative, using~\eqref{eq:sdof} to replace $\ddot z$, is
\begin{equation}
  \frac{\dd E}{\dd t} = \dot z\ddot z + \wn^2 z\dot z = \dot z\,(\ddot z + \wn^2 z)
  = \dot z\,(x - 2\zeta\wn\dot z) = \dot z\,x - 2\zeta\wn\dot z^2 .
\end{equation}
Take expectations: $\dd\E[E]/\dd t = \E[\dot z x] - 2\zeta\wn\E[\dot z^2]$. Two facts close this
equation.

\emph{Equipartition.} In a lightly damped oscillator the response is a sinusoid at $\wn$ with a
slowly varying amplitude $A$; over a cycle $\dot z^2$ and $\wn^2z^2$ have the same mean $\wn^2A^2/2$,
so $\E[\dot z^2] = \wn^2\E[z^2] = \E[E]$.

\emph{Input power.} The term $\E[\dot z(t)x(t)]$ is the power that the base injects. The velocity is a
weighted sum of past inputs, $\dot z(t) = \int_0^\infty h_v(u)x(t-u)\,\dd u$ ($h_v$ is the velocity
response to a unit impulse of $x$), so $\E[\dot z(t)x(t)] = \int_0^\infty
h_v(u)\E[x(t-u)x(t)]\,\dd u$. For a broadband input the correlation $\E[x(t-u)x(t)]$ vanishes after a
time much shorter than one period of the oscillator: the injected power depends only on the current
intensity of the input near $\fn$, not on the energy already stored. Its value follows from the
stationary case, where $\dd\E[E]/\dd t = 0$: then $\E[\dot zx] = 2\zeta\wn\E[E] = 2\zeta\wn^3\sigma_z^2$,
and with~\eqref{eq:miles}, $\E[\dot z x] = 2\zeta\wn^3\,G_0/(8\zeta\wn^3) = G_0/4$.

Let now the input intensity near $\fn$ change slowly, $G_0\to G(t)$, and write $s(t) = \E[z^2] =
\E[E]/\wn^2$ for the \emph{local variance} of the response. Dividing the energy balance by $\wn^2$:
\begin{equation}
  \frac{\dd s}{\dd t} = \frac{G(t)}{4\wn^2} - 2\zeta\wn\, s
  = \frac{1}{\tau}\bigl(s_\infty(t) - s(t)\bigr),
  \qquad s_\infty(t) = \frac{G(t)}{8\zeta\wn^3},
\end{equation}
where $s_\infty(t)$ is the stationary variance for the current input~\eqref{eq:miles} and
\begin{result}[Envelope law of the response]
\begin{equation}\label{eq:tau}
  \tau = \frac{1}{2\zeta\wn} = \frac{Q}{2\pi\fn},\qquad
  s(t) = \int_0^\infty \frac{1}{\tau}e^{-u/\tau}\, s_\infty(t-u)\,\dd u .
\end{equation}
\end{result}
The second form is the solution of the first-order equation (differentiate it to check). The
response variance is the stationary variance low-pass filtered with the time constant $\tau$. For a
non-white input the same law is applied with $s_\infty(t) = \int|H|^2G(f,t)\dd f$ computed from the
local spectrum $G(f,t)$.

\subsection{Consequence: fast modulations do not reach the response}\label{sec:papoulis}

Suppose the stationary variance $s_\infty(t)$ fluctuates about its mean with variance $V$ and an
exponential correlation $\E[\delta s_\infty(t)\delta s_\infty(t+u)] = Ve^{-|u|/T_c}$, $T_c$ being the
correlation time (the duration of a burst). In the frequency domain the fluctuation has the two-sided
spectrum $2VT_c/(1+\omega^2T_c^2)$ (its integral over $\omega/2\pi$ is $V$), and the filter
of~\eqref{eq:tau} has $|H_\tau(\omega)|^2 = 1/(1+\omega^2\tau^2)$. The variance of the filtered
fluctuation is
\begin{equation}
  \Var[s] = \int_{-\infty}^{\infty}\frac{\dd\omega}{2\pi}\,
  \frac{2VT_c}{(1+\omega^2T_c^2)(1+\omega^2\tau^2)} = \frac{2VT_c}{2\pi}\cdot\frac{\pi}{T_c+\tau}
  = V\,\frac{T_c}{T_c+\tau},
\end{equation}
where the integral $\int_{-\infty}^{\infty}\dd\omega/\bigl((1+\omega^2a^2)(1+\omega^2b^2)\bigr) =
\pi/(a+b)$ follows from the partial fractions $\frac{1}{(1+\omega^2a^2)(1+\omega^2b^2)} =
\frac{1}{a^2-b^2}\bigl(\frac{a^2}{1+\omega^2a^2}-\frac{b^2}{1+\omega^2b^2}\bigr)$ and
$\int\dd\omega/(1+\omega^2a^2) = \pi/a$: $\frac{1}{a^2-b^2}(\pi a - \pi b) = \pi/(a+b)$.
The mean of $s$ is not changed by the filter (its gain at $\omega=0$ is 1). Inserting in the last
form of~\eqref{eq:kmod}, $\kappa = 3\tavg{s^2}/\tavg{s}^2 = 3(1+\Var[s]/\tavg{s}^2)$, the kurtosis of
the response is
\begin{result}[Low-pass effect of the oscillator on the kurtosis]
\begin{equation}\label{eq:papoulis}
  \kappa_z - 3 = \bigl(\kappa_\infty - 3\bigr)\,\frac{T_c}{T_c+\tau},
\end{equation}
\end{result}
where $\kappa_\infty-3 = 3V/\tavg{s}^2$ is the excess kurtosis that the response would have if it
followed the input instantaneously. Bursts much shorter than $\tau$ leave a nearly Gaussian response
(this is Papoulis' rule); bursts much longer than $\tau$ are fully transmitted. Since $\tau$ is
inversely proportional to $\fn$, the same input can produce Gaussian responses at low natural
frequencies and leptokurtic ones at high natural frequencies. The kurtosis of the response,
$\kappa_z(\fn)$, is therefore a function of the natural frequency, and it is this function, not the
kurtosis of $x$, that governs the damage (Section~\ref{sec:damage}).

\subsection{The oscillator as the program computes it}\label{sec:digital}

The program computes $z$ from the samples $x_k$ with a second-order recursion,
\begin{equation}\label{eq:recursion}
  z_k = b_1x_k + b_2x_{k-1} + b_3x_{k-2} - a_1 z_{k-1} - a_2 z_{k-2},
\end{equation}
whose coefficients are those of the \emph{ramp-invariant} discretisation: they give the exact
response of~\eqref{eq:sdof} at the sampling instants when $x$ varies linearly between samples.
Appendix~\ref{app:smallwood} derives them. The frequency response of the recursion is obtained by
inserting $x_k = e^{\ii\omega k\Delta t}$, $z_k = H_d e^{\ii\omega k\Delta t}$:
\begin{equation}\label{eq:Hd}
  H_d(f) = \frac{b_1 + b_2 e^{-\ii\theta} + b_3 e^{-2\ii\theta}}{1 + a_1e^{-\ii\theta}+a_2e^{-2\ii\theta}},
  \qquad \theta = 2\pi f\Delta t .
\end{equation}
It equals $H$ when $\fn\ll f_s$ but not when $\fn$ approaches $f_s/2$. The ratio of the response
variances to a flat input, $\int|H_d|^2/\int|H|^2$, is for $Q=10$:
\begin{center}
\begin{tabular}{lccccc}
\toprule
$f_s/\fn$ & 25 & 8.8 & 6.25 & 4.4 & 3.3 \\
variance ratio & 0.99 & 0.92 & 0.85 & 0.72 & 0.56 \\
\bottomrule
\end{tabular}
\end{center}
The FDS of a record and of a test are both computed with~\eqref{eq:recursion}, so the comparison is
consistent; but a design based on the analogue~\eqref{eq:H2} overestimates the response at high
natural frequencies when $f_s$ is only a few times $f_{\max}$. The methods of version~2.0 design with
$H$ and let the closed loop (Section~\ref{sec:closedloop}) correct the difference; MMVSD designs with
$H_d$ directly (Section~\ref{sec:operators}).

\emph{Pre-filter.} When $f_s\le 10f_{\max}$ the program first passes $x$ through a fixed third-order
digital filter, forward and then backward in time (zero phase), as in the FDS computation of the
dissertation. Applied twice, the filter multiplies the PSD by $|H_{\mathrm{pre}}(f)|^4$. The same
filter is applied to records and tests, and $|H_{\mathrm{pre}}|^4$ multiplies $|H|^2$ in every
operator of the program, so that the designs stay consistent with the measurement.

\section{Fatigue damage and the Fatigue Damage Spectrum}\label{sec:damage}

\subsection{Basquin's law and Miner's rule}

A material loaded by cycles of constant stress amplitude $S$ fails after $N(S)$ cycles. In the range
of interest the S--N curve is a straight line in logarithmic scales (Basquin's law):
\begin{equation}
  N(S)\,S^b = C ,
\end{equation}
with the exponent $b$ (typically 3 to 12; 8 is a common default for structural details) and the
constant $C$. A cycle of amplitude $S$ consumes the fraction $1/N(S) = S^b/C$ of the life.
Miner's rule adds these fractions: a history containing $n_c$ cycles of amplitude $S_c$ has the damage
\begin{equation}\label{eq:miner}
  D = \sum_c \frac{n_c}{N(S_c)} = \frac1C\sum_c n_c S_c^b ,
\end{equation}
and failure is expected when $D$ reaches 1. With $S = Kz$ (Section~\ref{sec:oscillator}), the damage
of an oscillator is a sum over the cycles of its relative displacement:
\begin{equation}\label{eq:minerz}
  D = \frac{K^b}{C}\sum_c n_c\, A_c^b ,\qquad A_c = \text{amplitude (half range) of cycle } c .
\end{equation}
Only the ratio of two damages matters in test tailoring, so $K$ and $C$ are set to 1 by default.

\subsection{Rainflow counting}\label{sec:rainflow}

A random response does not consist of clean cycles; \emph{rainflow counting} decomposes it into
cycles in the way that is consistent with the stress--strain hysteresis of the material. The program
proceeds in two steps.
\begin{enumerate}[nosep]
\item \emph{Turning points.} Only the local maxima and minima of the sampled response matter; the
samples in between and plateaus are removed, leaving an alternating sequence $s_1, s_2, \dots$.
\item \emph{Four-point rule.} Take four consecutive turning points $s_{l}, s_{l+1}, s_{l+2},
s_{l+3}$ and the three ranges between them. If the inner range $|s_{l+2}-s_{l+1}|$ is not larger than
either outer range, the inner pair forms a closed cycle (a small loop inside a larger one): it is
counted with amplitude $|s_{l+2}-s_{l+1}|/2$ and the two points are removed. The rule is applied until
no more cycles close. The remaining sequence (the \emph{residue}) is counted as half cycles, one per
range, with weight $n_c = \tfrac12$.
\end{enumerate}
The program applies the four-point rule to all disjoint candidate pairs of the sequence at once, which
is fast in NumPy, and finishes long monotone decays with a sequential stack. Its damage sums are equal
to those of the sequential three-point algorithm of the ASTM E1049 standard (checked to $10^{-10}$ in
the tests).

\subsection{Fatigue Damage Spectrum and Maximum Response Spectrum}\label{sec:fds}

\begin{result}[FDS and MRS]
For a base acceleration $x$ of duration $T$, the Fatigue Damage Spectrum is the damage~\eqref{eq:minerz}
of the rainflow cycles of the response $z_{\fn}$ of each oscillator of the grid,
\begin{equation}\label{eq:fds}
  D(\fn) = r\,\frac{K^b}{C}\sum_{c\,\in\,\mathrm{rainflow}(z_{\fn})} n_c A_c^b ,
\end{equation}
and the Maximum Response Spectrum is the largest pseudo-acceleration,
\begin{equation}
  \mathrm{MRS}(\fn) = \wn^2\max_t|z_{\fn}(t)|.
\end{equation}
\end{result}
The factor $r$ is the number of times the signal is played: a record that represents a life
$\Tl$ longer than its duration $T$ has $r = \Tl/T$ (rounded to an integer when $\Tl\ge T$, i.e.\ the
record is considered repeated), and a test profile that is shorter than the test duration is played $r$
times (Section~\ref{sec:practice}). Damage adds over repetitions; the maximum does not, so the MRS has
no factor $r$. The MRS is a check of the peak stresses: a test whose MRS exceeds that of the life by
much may fail the component by overload rather than by fatigue.

\subsection{Rice's formulas for a Gaussian response}\label{sec:rice}

For a stationary Gaussian response the expected damage can be computed from the spectral moments.

\emph{Zero up-crossing rate.} The response crosses zero upwards in the interval $(t,t+\dd t)$ if
$z(t)\in(-\dot z\,\dd t, 0)$ with $\dot z>0$. The probability of this event is $\int_0^\infty
p(0,\dot z)\,\dot z\,\dd t\,\dd\dot z$, where $p(z,\dot z)$ is the joint density of $z$ and $\dot z$.
For a stationary process $z$ and $\dot z$ are uncorrelated, because $\E[z\dot z] =
\tfrac12\,\dd\E[z^2]/\dd t = 0$, and being jointly Gaussian they are independent:
$p(0,\dot z) = \frac{1}{\sqrt{2\pi}\sigma_z}\cdot\frac{1}{\sqrt{2\pi}\sigma_{\dot z}}
e^{-\dot z^2/(2\sigma_{\dot z}^2)}$. Dividing the probability by $\dd t$ and using
$\int_0^\infty \dot z\,e^{-\dot z^2/(2\sigma_{\dot z}^2)}\dd\dot z = \sigma_{\dot z}^2$,
\begin{equation}\label{eq:nu0}
  \nu_0^+ = \frac{\sigma_{\dot z}^2}{2\pi\sigma_z\sigma_{\dot z}} = \frac{\sigma_{\dot z}}{2\pi\sigma_z}
  = \sqrt{\frac{m_2}{m_0}} ,
\end{equation}
using $\sigma_{\dot z} = 2\pi\sqrt{m_2}$ and $\sigma_z = \sqrt{m_0}$ (Section~\ref{sec:respvar}).
The same reasoning applied to $\dot z$ (a maximum of $z$ is a down-crossing of zero by $\dot z$) gives
the rate of peaks $\nu_p = \sqrt{m_4/m_2}$. The ratio $\gamma = \nu_0^+/\nu_p = m_2/\sqrt{m_0m_4}$ is
the \emph{irregularity factor}: $\gamma=1$ means one peak per zero crossing (narrow band).

\emph{Amplitudes of a narrow-band response.}\label{sec:rayleigh} The response of a lightly damped
oscillator is narrow band: a sinusoid at $\fn$ whose amplitude changes slowly,
$z = A\cos(\wn t+\psi)$. Then $\dot z/\wn \approx -A\sin(\wn t+\psi)$ and the amplitude is
$A = \sqrt{z^2+(\dot z/\wn)^2}$. At a given time $z$ and $\dot z/\wn$ are independent zero-mean
Gaussian variables with the same variance $\sigma^2 = \sigma_z^2$ (equipartition,
Section~\ref{sec:envelope}). In polar coordinates $(A,\psi)$ their joint density
$\frac{1}{2\pi\sigma^2}e^{-A^2/(2\sigma^2)}$ times the Jacobian $A$, integrated over $\psi$, gives the
\emph{Rayleigh} density of the amplitude:
\begin{equation}
  p(A) = \frac{A}{\sigma^2}\,e^{-A^2/(2\sigma^2)},\qquad A\ge0 .
\end{equation}
Its moment of order $b$ follows with the substitution $u = A^2/(2\sigma^2)$, $A = \sqrt{2\sigma^2u}$,
$\dd u = A\,\dd A/\sigma^2$:
\begin{equation}\label{eq:rayleighmoment}
  \E[A^b] = \int_0^\infty A^b\,\frac{A}{\sigma^2}e^{-A^2/(2\sigma^2)}\,\dd A
  = \int_0^\infty (2\sigma^2u)^{b/2}e^{-u}\,\dd u = (2\sigma^2)^{b/2}\,\Gamma\!\left(1+\tfrac b2\right).
\end{equation}

\emph{Expected damage.} A narrow-band response has one cycle per period, $\nu_0^+T \approx \fn T$
cycles in the duration $T$, each with a Rayleigh amplitude. By~\eqref{eq:minerz} the expected damage is
the number of cycles times $\E[A^b]$:
\begin{result}[Gaussian narrow-band damage (Rice--Lalanne)]
\begin{equation}\label{eq:dgauss}
  D_G(\fn) = \frac{K^b}{C}\,\fn\,T\,\bigl(2\sigma_z^2\bigr)^{b/2}\,\Gamma\!\left(1+\tfrac b2\right),
  \qquad \sigma_z^2 = \sum_j W_{ij}G_j .
\end{equation}
\end{result}
Solving for the variance gives the response variance that produces a prescribed damage $D$ in the
duration $T$:
\begin{equation}\label{eq:targetvar}
  \sigma_z^2 = \frac12\left(\frac{D\,C}{K^b\,\fn\,T\,\Gamma(1+b/2)}\right)^{2/b}.
\end{equation}
For wide-band responses (for instance the quasi-static response of a high-frequency oscillator to a
strong low-frequency input) the cycles are not one per period and their amplitudes are not Rayleigh
distributed. Dirlik's empirical distribution of the rainflow ranges (Appendix~\ref{app:dirlik})
covers this case; it reduces to~\eqref{eq:dgauss} in the narrow-band limit $\gamma\to1$. MMVSD uses
it (Section~\ref{sec:damagemodel}).

\subsection{Rainflow on samples: peaks between samples}\label{sec:sampling}

Rainflow counting works on samples. The true maximum of a cycle generally falls between two samples;
the largest sample of the half cycle is the one nearest to it. For a response of frequency $f$ the
nearest sample lies within $\pm\Delta t/2$ of the true peak, at a phase offset $\vartheta = 2\pi f
\delta$ uniformly distributed in $(-\pi f/f_s, \pi f/f_s)$, and the counted peak is $A\cos\vartheta$
instead of $A$. A range is measured between a maximum and the following minimum, with independent
offsets $\vartheta_1,\vartheta_2$; its amplitude is $A(\cos\vartheta_1+\cos\vartheta_2)/2$, and its
damage is reduced on average by the factor
\begin{result}[Sampling factor]
\begin{equation}\label{eq:sampling}
  c_b(\rho) = \E\!\left[\left(\frac{\cos\vartheta_1+\cos\vartheta_2}{2}\right)^{\!b}\right]
  = \frac{1}{(\pi\rho)^2}\int_0^{\pi\rho}\!\!\int_0^{\pi\rho}
  \left(\frac{\cos\vartheta_1+\cos\vartheta_2}{2}\right)^{\!b}\dd\vartheta_1\,\dd\vartheta_2 ,
  \qquad \rho = \frac{f}{f_s},
\end{equation}
\end{result}
(the integrand is even in each angle, so the average over $(-\pi\rho,\pi\rho)$ equals the average
over $(0,\pi\rho)$). For $b=8$: $c_b = 0.98$ at $\rho = 0.04$, $0.86$ at $\rho = 0.11$ and $0.54$ at
$\rho=0.23$. The program evaluates~\eqref{eq:sampling} by Gauss--Legendre quadrature on a table of
$\rho$. Together with the digital response of Section~\ref{sec:digital} this explains why a Gaussian
design underestimates the measured damage loss at high natural frequencies when $f_s$ is close to
$2.5f_{\max}$; MMVSD includes both effects in its model.

\subsection{Why the kurtosis of the response matters}\label{sec:kzdamage}

Let the response be Gaussian given a slowly varying local variance $s(t)$ (Section~\ref{sec:modulated}).
Over a stretch of time where $s$ is constant, the damage rate is the narrow-band
rate~\eqref{eq:dgauss} with $\sigma_z^2 = s$. Averaging over time,
\begin{equation}
  D = \frac{K^b}{C}\fn T\,\Gamma\!\left(1+\tfrac b2\right) 2^{b/2}\,\tavg{s^{b/2}}
  = D_G\bigl(\tavg{s}\bigr)\;\frac{\tavg{s^{b/2}}}{\tavg{s}^{b/2}} ,
\end{equation}
where $D_G(\tavg s)$ is the Gaussian damage of a stationary response with the same variance.
The factor $\tavg{s^{b/2}}/\tavg{s}^{b/2}$ is at least 1 (Jensen's inequality, since $s^{b/2}$ is
convex for $b\ge2$). For $b=4$ it equals $\tavg{s^2}/\tavg{s}^2 = \kappa_z/3$ by~\eqref{eq:kmod}; for
larger $b$ it depends on higher moments of $s$ but grows with $\kappa_z$. Two signals with the same PSD
can therefore produce very different damage if their responses have different kurtosis, and a test
that reproduces the FDS with a different $\kappa_z(\fn)$ reproduces it with a different mix of cycles
(many moderate cycles instead of a few large ones). This is the reason for prescribing $\kappa_z$
together with the FDS (Part~III).

\part*{Part II --- Mission synthesis}
\addcontentsline{toc}{part}{Part II --- Mission synthesis}

\section{Synthesis of a profile with a prescribed FDS}\label{sec:synthesis}

\subsection{The tailoring problem}

A target is described by its FDS $D_t(\fn)$ over the life $\Tl$ (Section~\ref{sec:fds}), its MRS,
its PSD and its statistics. A specification (a PSD or a swept sine) or a mission
(Section~\ref{sec:missions}) is converted into the same description. The test must last $T\le\Tl$ and
reproduce $D_t$ on every node. If the life is compressed by the factor $\Tl/T$, the damage per unit
time must increase by the same factor; with the narrow-band law~\eqref{eq:dgauss}, where $D\propto
T\sigma_z^b$, this means multiplying the response amplitudes, hence the test amplitudes, by
$(\Tl/T)^{1/b}$. For $b=8$ a ten times shorter test is only $10^{1/8} = 1.33$ times more severe in
amplitude; this mild growth is what makes accelerated testing possible.

\subsection{Gaussian design: from the FDS to a PSD}\label{sec:inversion}

The standard mission synthesis looks for a stationary Gaussian profile. Equation~\eqref{eq:targetvar}
with the test duration $T$ gives, at every node, the response variance $\sigma^2_{t,i}$ that produces
the target damage. Equation~\eqref{eq:W} relates the response variances to the PSD values, so the
design is the linear problem
\begin{equation}\label{eq:WG}
  \sum_j W_{ij}G_j = \sigma^2_{t,i}, \qquad G_j\ge0 .
\end{equation}
All entries of $W$ are positive and each row is dominated by the columns near its own node (the
oscillator responds mainly to its band). The program solves~\eqref{eq:WG} by the multiplicative
fixed point
\begin{equation}\label{eq:multfix}
  G_j \leftarrow G_j\,\frac{\sigma^2_{t,j}}{(WG)_j},
\end{equation}
started from $G_j = \sigma^2_{t,j}/W_{jj}$. If the response of node $j$ is too small,
$(WG)_j<\sigma^2_{t,j}$ and $G_j$ grows; if too large, it shrinks; positivity is preserved because
only positive factors are applied. The iteration stops when $\max_j|\sigma^2_{t,j}/(WG)_j - 1| <
10^{-4}$; when it stalls, a non-negative least-squares solution of the scaled system is used instead.
Because neighbouring oscillators overlap, the exact solution may set some nodes to zero and raise
their neighbours; the synthesis keeps a tiny floor on such lines so that later multiplicative
corrections can revive them.

\subsection{Synthesis of the time history}

The PSD on the FFT grid of the test (line spacing $1/T$) is the linear interpolation of the nodal
values, held constant over $3/Q$ relative bandwidth beyond the first and last node. The time history
is the random-phase multisine of Section~\ref{sec:multisine} with these line powers. The phases are
drawn once with the random seed and then frozen.

\subsection{The closed loop}\label{sec:closedloop}

The design uses the Gaussian formula~\eqref{eq:dgauss}, while the test will be judged by the rainflow
FDS~\eqref{eq:fds} of the realised signal; the two differ (wide-band responses, the digital oscillator,
the sampling of the peaks, and the scatter of a single realisation). Every FDS method of the package
therefore ends with a \emph{closed loop}: the rainflow FDS $D_i$ of the realisation is measured and the
spectrum is corrected until it matches the target.

\emph{Correction step.} If the amplitude of the signal in the band of node $i$ is multiplied by $g_i$,
the response amplitudes of oscillator $i$ are multiplied by $g_i$ and, by~\eqref{eq:minerz}, its damage
by $g_i^b$. The gain that would correct an isolated node is therefore $(D_{t,i}/D_i)^{1/b}$. The
program uses
\begin{equation}\label{eq:gain}
  g_i = \mathrm{clip}\!\left(\left(\frac{D_{t,i}}{D_i}\right)^{w/b},\ \tfrac{1}{1.5},\ 1.5\right),
  \qquad G(f)\leftarrow G(f)\,g(f)^2,
\end{equation}
where $g(f)$ is the linear interpolation of the nodal gains, the phases stay frozen and, for the
modulated methods, the envelope is kept. The relaxation $w$ starts at 1 and grows by 25\,\% after each
improving step up to 1.6 (over-relaxation: because neighbouring oscillators share their bands, a
plain step corrects only part of the error). A step that makes the result worse is undone and $w$ is
halved.

\emph{Acceptance score.} A step is judged by $\tfrac12\max_i e_i + \tfrac12 q_{90}(e)$, where
$e_i = |D_i/D_{t,i}-1|$ and $q_{90}$ is the 90th percentile over the nodes: one stubborn oscillator
must not block the improvement of all the others. The loop stops when the maximum error is below the
tolerance (5\,\% by default), when the number of corrections is exhausted, or when $w<1/16$; the best
realisation found is returned.

\subsection{Long tests}\label{sec:repeats}

A test longer than a maximum number of samples (20 million by default; 4 million for MMVSD) is
synthesised as a shorter profile of duration $T/r$ that is played $r$ times; the target of the
profile is $D_t/r$. Damage adds over repetitions (Section~\ref{sec:fds}), so the whole test meets
$D_t$.

\section{The methods of version 2.0}\label{sec:methods}

All methods of this section are those of the PhD dissertation, written in a common form. The
FDS-controlled methods (\code{gaussian}, \code{kfds}, \code{rf}, \code{pf}, \code{psf}, \code{sor})
end with the closed loop of Section~\ref{sec:closedloop}.

\subsection{kFDS: prescribed FDS and kurtosis}\label{sec:kfds}

\emph{Form of the signal.} kFDS multiplies a stationary Gaussian carrier $g$, designed as in
Section~\ref{sec:inversion}, by a slowly varying envelope $e(t)$,~\eqref{eq:envcarrier}. The envelope
is piecewise constant over windows of length $L_w$, with linear ramps of $L_w/3$ between the levels, so
that the waveform has no jumps.

\emph{The damage is kept by the normalisation $\tavg{e^b} = 1$.} If a window is long compared with the
response time of the oscillators, the response in the window is the carrier's response multiplied by
$e_w$, its cycles are those of the carrier scaled by $e_w$, and by~\eqref{eq:minerz} the damage in the
window is $e_w^b$ times the carrier's damage in the window. If the carrier's damage is spread evenly
over time, summing over the windows gives
\begin{equation}
  D(e\,g) = \sum_w e_w^b\,D_{g,w} = \sum_w e_w^b\,\frac{T_w}{T}D_g = \tavg{e^b}\,D_g ,
\end{equation}
so any envelope with $\tavg{e^b}=1$ keeps the FDS of the carrier. (In the dissertation the same
condition appears as weights $v_j$ of blocks, $e_j = (n v_j)^{1/b}$.)

\emph{The kurtosis is set by the contrast of the envelope.} By~\eqref{eq:kmod} the kurtosis of the
product is $3\tavg{e^4}/\tavg{e^2}^2$. The envelope is $e_w = \exp(\alpha u_w)$ normalised to
$\tavg{e^b}=1$, where the relative log-levels $u_w\le0$ are drawn
\begin{itemize}[nosep]
\item from the distribution of the window RMS of the reference (\code{pattern = "reference"},
default when a record is available): the logarithm of the RMS of every window of the record is
computed, and $n_w$ values are drawn from their life-weighted distribution by quantile resampling, in
random order, so that the test has bursts with the statistics of the field data;
\item or as random bursts (\code{"bursts"}): 10 to 20\,\% of the windows at a high level.
\end{itemize}
$\alpha = 0$ gives a constant envelope (kurtosis 3); the kurtosis grows with $\alpha$. The program
finds the $\alpha$ that gives the target kurtosis \emph{on the realised signal} $e\,g$ with a
bracketing root finder (Brent's method). The envelope is frozen after the first two iterations of the
closed loop, which then corrects only the carrier spectrum.

\emph{Window length.} The normalisation argument requires windows long compared with the response
time of the slowest oscillator. The amplitude of a free oscillator decays as $e^{-\zeta\wn t}$, with the
decay time $1/(\zeta\wn) = Q/(\pi\fn) = 2\tau$; the window is at least $2Q/(\pi f_{\min})$. Within this
bound the default follows the rule of the dissertation: a fraction of the reference duration (1/150 to
1/400 depending on $f_s$), with at least 300 samples.

\subsection{RF: prescribed FDS and RMS}

RF uses the same envelope with the RMS as the target. With $\tavg{e^b}=1$ the RMS of the product is
$\sigma_g\tavg{e^2}^{1/2}$. For $b>2$ the power-mean inequality gives $\tavg{e^2}^{1/2}\le
\tavg{e^b}^{1/b} = 1$, with equality only for a constant envelope; increasing the contrast lowers the
RMS. RF can therefore only lower the RMS below that of the Gaussian profile with the same FDS; the
program says so when a higher RMS is requested.

\subsection{PF and PSF: prescribed FDS and PSD}\label{sec:pf}

PF builds the test from two independent stationary Gaussian carriers, a high-energy one $x_B$ (PSD
$G_B$) for the bursts and a low-energy one $x_L$ (PSD $G_L$) for the rest, combined with a smooth
switching mask $m(t)\in[0,1]$ that is 1 during the bursts:
\begin{equation}
  x(t) = \sqrt{m(t)}\,x_B(t) + \sqrt{1-m(t)}\,x_L(t).
\end{equation}
Because the carriers are independent, the local PSD is $mG_B + (1-m)G_L$, and the time-averaged PSD is
$\tavg{m}G_B + (1-\tavg m)G_L = (t_BG_B + t_LG_L)/T$, with $t_B = \tavg mT$ the total burst time and
$t_L = T-t_B$. Requiring the reference PSD gives
\begin{equation}
  G_L = \frac{T\,G_{\mathrm{ref}} - t_B\,G_B}{t_L}.
\end{equation}
The damage is the sum of the damages of the two parts. The burst PSD is found by a fixed point:
start with $D_L = D_t/20$; design $G_B$ for $D_B = D_t-D_L$ over $t_B$ (Section~\ref{sec:inversion});
compute $G_L$ from the PSD condition and its Gaussian damage $D_L$ over $t_L$; repeat. The burst time
starts at $T/20$ and is halved while $G_L$ would be negative over more than 10\,\% of the band, but
two bursts of at least three decay times of the lowest oscillator are kept when possible. When
PSD, RMS and FDS cannot all be met (for instance at high acceleration), the program reports it.

PSF prescribes only the shape of the PSD: PF is run with the damage that the reference produces in
the test duration, $D_t\,T/\Tl$, and the result is multiplied by $(\Tl/T)^{1/b}$, which multiplies
its damage by $\Tl/T$ and its PSD by $(\Tl/T)^{2/b}$ without changing the shape.

\subsection{Kurtosis control without damage: MLV, VSD, PM}

These methods prescribe the PSD and the kurtosis of the signal, not its damage. MLV
(multi-level variance) is the modulated Gaussian~\eqref{eq:envcarrier} normalised to $\tavg{e^2}=1$
(this keeps PSD and RMS) with the contrast solved for the kurtosis. VSD (variable spectral density)
splits the band into log-spaced sub-bands, each with its own carrier and independent burst envelope,
so that the shape of the PSD changes with time while its time average is the reference PSD. PM
(phase manipulation) keeps Gaussian amplitudes and partially aligns the phases of the harmonics in a
fraction of the blocks, which creates peaks; the degree of alignment is solved for the kurtosis. The
benchmark of the program shows that, over the same duration, their FDS can differ from the
reference's by a factor of 3 or more: kurtosis alone does not control the damage.

\subsection{MRS control}

The largest value of a narrow-band Gaussian response in the duration $T$ can be estimated from Rice's
crossing rate. Repeating the derivation of~\eqref{eq:nu0} for the level $z=u$ instead of 0, the
joint density at $z=u$ contains the extra factor $e^{-u^2/(2\sigma_z^2)}$, so the rate of up-crossings
of $u$ is $\nu_0^+e^{-u^2/(2\sigma_z^2)}$. The level crossed on average once in $T$ satisfies
$\nu_0^+Te^{-u^2/(2\sigma_z^2)} = 1$, i.e.
\begin{equation}
  u = \sigma_z\sqrt{2\ln(\nu_0^+T)} \approx \sigma_z\sqrt{2\ln(\fn T)} ,
  \qquad \mathrm{MRS}_G(\fn) = \wn^2\sigma_z\sqrt{2\ln(\fn T)} .
\end{equation}
A target MRS therefore prescribes $\sigma_z$ at every node, the PSD follows from~\eqref{eq:WG}, and
the closed loop corrects the measured MRS with the exponent 1 instead of $1/b$ in~\eqref{eq:gain}
(the MRS is proportional to the amplitude).

\subsection{Sine on random}\label{sec:sor2}

Rotating machines add tonal lines to the random vibration. The method \code{sor} detects the lines
in the reference PSD: a peak standing more than a prominence threshold (10\,dB by default) above a
running median of the PSD is a line, and its amplitude follows from the power in the peak,
$A = \sqrt{2P}$ (a sine of amplitude $A$ has the power $A^2/2$). A stationary sine of amplitude $A$
produces in an oscillator a sinusoidal response of amplitude $|H|A$ at the line frequency, i.e.\
the damage $\frac{K^b}{C}f_0 T(|H|A)^b$ in the time $T$. To keep the damage of the life $\Tl$ in the
test duration $T$, the test amplitude $A_T$ must satisfy $A_T^bT = A^b\Tl$:
\begin{equation}\label{eq:sorgain}
  A_T = A\left(\frac{\Tl}{T}\right)^{1/b}.
\end{equation}
The random part is designed for the remaining damage and corrected in closed loop with the lines kept
fixed; nodes where the lines alone exceed the target are reported and left out of the loop criterion.
Version~2.1 adds the detection of lines that are present only in some blocks of a record
(Section~\ref{sec:lines}).

\subsection{Missions}\label{sec:missions}

A mission is a tree whose leaves are situations (a record or a specification with the life it
represents) and whose inner nodes are groups. In a \emph{series} group the component lives through
all children one after the other: the FDS are added, the lives are added, the PSD is the
life-weighted mean, the MRS is the maximum. In a \emph{parallel} group it lives through only one of
them, not known in advance: the FDS and PSD are enveloped (maximum over the children) and the life,
RMS and kurtosis are those of the most damaging child.

The statistics of a series group follow from the moments of a mixture. If a fraction $w_c$ of the
life is spent in situation $c$, where the signal has the moments $\mu_{2,c}$ and $\mu_{4,c}$, the
moments over the whole life are the weighted averages $\mu_2 = \sum_cw_c\mu_{2,c}$ and $\mu_4 =
\sum_cw_c\mu_{4,c}$, and the kurtosis is $\mu_4/\mu_2^2$. Version~2.1 applies the same rule to the
response of every oscillator: with $v_c(\fn)$ and $\kappa_c(\fn)$ the variance and kurtosis of the
response in situation $c$, $\mu_{4,c} = \kappa_cv_c^2$ and
\begin{result}[Response kurtosis of a series mission]
\begin{equation}\label{eq:kzmix}
  v(\fn) = \sum_c w_c v_c(\fn),\qquad
  \kappa_z(\fn) = \frac{\sum_c w_c\,\kappa_c(\fn)\,v_c(\fn)^2}{v(\fn)^2}.
\end{equation}
\end{result}
A PSD specification has $\kappa_c = 3$ (it describes a Gaussian process); a swept sine is measured on
its generated signal.

\part*{Part III --- MMVSD}
\addcontentsline{toc}{part}{Part III --- MMVSD}

\section{The MMVSD model}\label{sec:mmvsdmodel}

\subsection{Why one envelope is not enough}

In kFDS the same envelope multiplies all frequencies. When the windows are long compared with the
response times, every oscillator follows the envelope and, by~\eqref{eq:kmod}, every response has the
same kurtosis $\kappa_z = 3\tavg{e^4}/\tavg{e^2}^2$, whatever $\fn$. Measured vibration rarely
behaves like this. A vehicle combines a road input whose level changes every few seconds (leptokurtic
responses at low $\fn$), an engine that excites a middle band stationarily (Gaussian responses there)
and impacts that excite high frequencies in short bursts (strongly leptokurtic responses at high
$\fn$). The kurtosis of the responses is a function $\kappa_z(\fn)$, and a test that reproduces the FDS
with a different function reproduces the damage with a different mix of cycles
(Section~\ref{sec:kzdamage}).

MMVSD (Moment-Matched Variable Spectral Density) prescribes both functions: the FDS $D_t(\fn)$ and the
kurtosis of the responses $\kappa_t(\fn)$, measured on the reference (or combined over a mission
with~\eqref{eq:kzmix}). It extends the idea of a variable spectral density, a PSD whose shape changes
with time, to a continuous random field designed so that the two families of moments of the
responses, their damage and their kurtosis, match the target.

\subsection{An evolutionary spectrum}\label{sec:evol}

The test signal is described by a \emph{local spectrum} that changes with time:
\begin{result}[MMVSD signal model]
\begin{equation}\label{eq:evol}
  G(f,t) = P(f)\,m(f,t),\qquad m(f,t)\ge0,\qquad \tavg{m(f,\cdot)} = 1 \ \text{for every } f,
\end{equation}
and, given the field $m$, the signal is a Gaussian process with this local spectrum.
\end{result}
The time average of the local spectrum is $\tavg{G(f,\cdot)} = P(f)\tavg{m(f,\cdot)} = P(f)$: the PSD
of the test is exactly $P$, whatever the modulation. The \emph{modulation field} $m$ says when and in
which bands the energy is concentrated; its distribution sets how leptokurtic the responses are.

$P$ is parametrised as a smooth correction of a \emph{shape prior} $\Pi(f)$, the PSD of the reference
at its full resolution:
\begin{equation}\label{eq:prior}
  P(f) = \Pi(f)\,p(f),\qquad p(f) = \sum_j p_j\,\phi_j(f),
\end{equation}
where $\phi_j$ are the hat functions of the nodes (Section~\ref{sec:lalanne}) extended flat over $3/Q$
beyond the first and last node. The peaks, notches and band edges of the reference are thus kept at
full resolution, and the unknowns are the $n$ nodal factors $p_j$. Outside the band of the grid the
test has no energy. The modulation is likewise described by its nodal values $m_j(t)$.

\subsection{Local moments of the responses}\label{sec:operators}

For a stationary input the spectral moments of the response of oscillator $i$
are~\eqref{eq:moments}. For the local spectrum~\eqref{eq:evol}, with $P$ and $m$ written through their
nodal values, the moments that the response would have if the local spectrum at time $t$ lasted for
ever are
\begin{equation}\label{eq:Sinf}
  S^{\infty}_{k,i}(t) = \sum_j W^{(k)}_{ij}\,p_j\,m_j(t),
  \qquad
  W^{(k)}_{ij} = \int f^k\,|H_{d,i}(f)|^2\,|H_{\mathrm{pre}}(f)|^4\,\Pi(f)\,\phi_j(f)\,\dd f,
\end{equation}
for $k = 0,1,2,4$. The \emph{moment operators} $W^{(k)}$ generalise Lalanne's operator ($W^{(0)}$ with
$\Pi=1$, the analogue response and no edge extension is the matrix $W$ of~\eqref{eq:W}). MMVSD uses
the response $H_d$ of the digital oscillator~\eqref{eq:Hd} and the pre-filter of
Section~\ref{sec:digital}, so that the model describes exactly the quantities that the measurement
computes. The integrals are evaluated by the trapezoidal rule with 128 intervals per segment between
nodes. The model assumes that within the band of an oscillator the local spectrum is the
interpolation of its nodal values, i.e.\ that the modulation varies slowly along the frequency axis
compared with the bandwidth $\fn/Q$.

The response does not follow the local spectrum instantly but with the envelope law~\eqref{eq:tau}.
Applying the low-pass filter $F_i$ of time constant $\tau_i = Q/(2\pi f_{\mathrm{n},i})$ to every
moment gives the \emph{local moments}
\begin{result}[Local moments of the responses]
\begin{equation}\label{eq:S}
  S_{k,i}(t) = F_i\!\left[\,\sum_j W^{(k)}_{ij}\,p_j\,m_j(\cdot)\right]\!(t),\qquad
  F_i[y](t) = \int_0^\infty \frac{1}{\tau_i}e^{-u/\tau_i}\,y(t-u)\,\dd u .
\end{equation}
\end{result}
$s_i(t) = S_{0,i}(t)$ is the local variance of the response of oscillator $i$. On the discrete time
grid of the model (step $\delta$, see Section~\ref{sec:frames}) the filter becomes the recursion
$y^F_l = a_iy^F_{l-1} + (1-a_i)y_l$ with $a_i = e^{-\delta/\tau_i}$, whose gain at zero frequency is
$(1-a_i)/(1-a_i) = 1$. The test profile is periodic (it may be played several times), so the
recursion is applied circularly, by multiplying the discrete Fourier transform of $y$ by
$(1-a_i)/(1-a_ie^{-2\pi\ii q/N_t})$ at the frequency index $q$ of the $N_t$ points of the grid.

\subsection{Kurtosis of the responses}\label{sec:kurtmodel}

Given $m$, the response of oscillator $i$ is Gaussian with the local variance $s_i(t)$. By the last
form of~\eqref{eq:kmod},
\begin{result}[Response kurtosis in the model]
\begin{equation}\label{eq:kappamodel}
  \kappa_i = 3\,\frac{\tavg{s_i^2}}{\tavg{s_i}^2}.
\end{equation}
\end{result}
Note that $\tavg{s_i} = \sum_jW^{(0)}_{ij}p_j\tavg{m_j} = \sum_jW^{(0)}_{ij}p_j$ does not depend on
the modulation (the filter keeps the mean and $\tavg{m_j}=1$): the modulation changes the kurtosis
and the damage, never the mean response variance.

\subsection{Damage of the responses}\label{sec:damagemodel}

Over a stretch of time short compared with the variations of $s_i$ but containing many cycles, the
response is a stationary Gaussian process with the moments $S_{k,i}(t)$. Its damage rate is the
expected damage per second of such a process: Dirlik's rate (Appendix~\ref{app:dirlik}), a function
$d_{\mathrm{Dir}}(S_0,S_1,S_2,S_4)$ that reduces to the narrow-band rate $\nu_0^+(2S_0)^{b/2}
\Gamma(1+b/2)$ for narrow bands, multiplied by the sampling factor~\eqref{eq:sampling} evaluated at
the mean crossing rate $\nu_0^+ = \sqrt{S_2/S_0}$ (the cycles that dominate the damage are those of
the zero crossings):
\begin{result}[Damage in the model]
\begin{equation}\label{eq:Dmodel}
  D_i = \frac{K^b}{C}\,T\,\Bigl\langle d\bigl(S_{0,i},S_{1,i},S_{2,i},S_{4,i}\bigr)\Bigr\rangle ,
  \qquad
  d = d_{\mathrm{Dir}}(S_0,S_1,S_2,S_4)\;c_b\!\left(\frac{\sqrt{S_2/S_0}}{f_s}\right).
\end{equation}
\end{result}
For a stationary Gaussian signal ($m\equiv1$) and a narrow band this is Lalanne's
formula~\eqref{eq:dgauss}, corrected for the digital oscillator and the sampling. The program also
offers the plain narrow-band (Rice) rate instead of Dirlik's.

How well does the model describe a realisation? For a stationary Gaussian profile synthesised by the
program at $f_s = 500$\,Hz with nodes up to 150\,Hz ($Q=10$, $b=8$), the measured response variances
agree with $S_0$ within 2--7\,\% at all nodes, and the ratio of the rainflow damage to~\eqref{eq:Dmodel}
averages about 0.9 (the remaining difference is the scatter of the largest cycles of a 60\,s
realisation, whose median is below its mean). With the analogue operators and without the sampling
factor, the same ratio drops to 0.15 at 113\,Hz.

\subsection{The modulation field}\label{sec:field}

The field must have, at every node, a prescribed depth of modulation, and between nodes and over time
the correlations that describe the bursts of the reference: how long a burst lasts (\emph{burst
duration} $T_c$, possibly different at every node) and how wide a band it covers (\emph{burst width}
$B$, in octaves). It is built in three steps on a grid of $K$ time frames (spacing $h_t$, the
\emph{hop}).

\emph{1. A latent Gaussian field.} Draw independent standard normal numbers $\xi_{j\ell}$ for every
node $j$ and frame $\ell$. Smooth every row in time with a Gaussian kernel of standard deviation
$T_{c,j}/h_t$ frames (circularly, by multiplying its discrete Fourier transform by
$\exp(-2\pi^2\varsigma_j^2\nu^2)$, $\nu$ in cycles per frame, $\varsigma_j = T_{c,j}/h_t$), standardise
each row to unit variance, then mix the rows with the Gaussian kernel
$\exp\bigl(-\tfrac12(\log_2f_{\mathrm{n},j}-\log_2f_{\mathrm{n},j'})^2/B^2\bigr)$ along the logarithmic
frequency axis, and standardise again. The result $Z_{j\ell}$ is a Gaussian field with zero mean, unit
variance at every node, and the prescribed correlations. A Gaussian kernel of standard deviation
$\varsigma$ applied to white noise gives the correlation $e^{-\lambda^2/(4\varsigma^2)}$ at the lag
$\lambda$ (the autocorrelation of the kernel is a Gaussian of variance $2\varsigma^2$); the correlation
falls to $1/e$ at $\lambda = 2\varsigma$. The settings $T_c$ and $B$ are therefore half the $1/e$
widths of the correlations.

\emph{2. Ranks.} In every row, replace the values by their ranks $r_{j\ell}\in\{0,\dots,K-1\}$ over
time. The ranks keep the information of when the field is high or low, which is all that the copula
(the dependence structure) needs.

\emph{3. Marginal.} The Gamma distribution of shape $a>0$ and scale $1/a$ has the density
$p(m) = a^a m^{a-1}e^{-am}/\Gamma(a)$ for $m>0$. With the substitution $u = am$, its moments are
$\E[m^k] = \frac{1}{a^k\Gamma(a)}\int_0^\infty u^{a+k-1}e^{-u}\dd u = \frac{\Gamma(a+k)}{a^k\Gamma(a)}$;
since $\Gamma(s+1) = s\Gamma(s)$ (integrate by parts), $\E[m] = \frac{a\Gamma(a)}{a\Gamma(a)} = 1$ and
$\E[m^2] = \frac{(a+1)a}{a^2} = 1+\frac1a$, so the variance is $1/a$. Map the ranks to the quantiles of
this distribution with $a_j = 3/e_j$ (mean 1, variance $e_j/3$):
\begin{equation}\label{eq:gammafield}
  m_{j\ell} = F^{-1}_{a_j}\!\left(\frac{r_{j\ell}+\tfrac12}{K}\right),
\end{equation}
and divide every row by its mean so that $\tavg{m_j} = 1$ exactly. Because every row uses the same set
of $K$ quantiles, the sample moments of $m_j$ depend only on $e_j$: there is no sampling noise in
$\tavg{m_j^2}$ or $\tavg{m_j^{b/2}}$ (a \emph{stratified} copula). The quantiles of the Gamma
distribution are tabulated for 200 shapes and interpolated, which makes the field cheap to rebuild
during the design.

\emph{Meaning of the depth $e_j$.} A Gaussian signal modulated directly by $m$ (local variance
proportional to $m$) has, by~\eqref{eq:kmod}, the kurtosis $3\tavg{m^2}/\tavg{m}^2 =
3(1+\Var m) = 3 + e$; $e$ is the \emph{excess kurtosis of the modulation}. The kurtosis of a
response is lower when the bursts are short compared with $\tau_i$, by~\eqref{eq:papoulis}; the design
(Section~\ref{sec:design}) finds the $e_j$ that give the target $\kappa_t$ through the full
model~\eqref{eq:S}--\eqref{eq:kappamodel}, so this reduction is accounted for. The Gamma family is
positive, has one shape parameter, contains the exponential distribution ($e=3$) and tends to the
constant 1 as $e\to0$. The depth is capped (30 by default) to avoid unrealistic peaks.

\subsection{Burst width and duration from the record}\label{sec:structure}

The program estimates $B$ and $T_{c,j}$ from the reference. The record is cut into frames of 0.5\,s
(hop 0.25\,s, Hann window, i.e.\ the taper $\sin^2(\pi n/L)$ over the $L$ samples of a frame); in every
frame the power in the 1/3-octave band around every node is computed, and its logarithm
$\beta_{j\ell}$ is taken. Even for a stationary signal these values scatter: the power of a band
of $n_l$ lines is a sum of $2n_l$ squared Gaussian variables (the real and imaginary parts of each
line), i.e.\ a \emph{chi-square} variable with $\nu = 2n_l$ degrees of freedom up to a scale factor.
Half of such a variable, $Y$, has the Gamma distribution of shape $k = \nu/2$ and scale 1 (density
$y^{k-1}e^{-y}/\Gamma(k)$), whose moments are $\E[Y^s] = \int_0^\infty y^{s+k-1}e^{-y}\dd y/\Gamma(k) =
\Gamma(k+s)/\Gamma(k)$. The logarithm of $\E[e^{s\ln Y}] = \E[Y^s]$, the cumulant generating function
of $\ln Y$, is therefore $\ln\Gamma(k+s)-\ln\Gamma(k)$; its second derivative at $s=0$ is the
variance of $\ln Y$, the \emph{trigamma function} $\psi'(k) = \frac{\dd^2}{\dd k^2}\ln\Gamma(k)$ at
$k=\nu/2$ (a scale factor does not change the variance of a logarithm). This noise variance is subtracted from the variance of each band before the
correlations are normalised (the effective $\nu$ of a band with $n_l$ lines is taken as $2n_l/1.5$,
since the Hann window correlates neighbouring lines). Bands whose remaining variance is negligible
do not fluctuate and are left out. Then:
\begin{itemize}[nosep]
\item the time correlation of each band is its normalised autocorrelation; the lag where it falls
below $1/e$, halved, is $T_{c,j}$ (smoothed along the nodes and interpolated to all of them);
\item the frequency correlation is the average correlation between bands as a function of their
distance in octaves; the distance where it falls below $1/e$, halved, is $B$ (between 1/12 and 3
octaves).
\end{itemize}
The burst duration is finally kept above $\max(2\tau_i, 0.25\,\mathrm{s})$: shorter bursts would not
reach the responses (Section~\ref{sec:papoulis}) and would only require deeper modulation.

\subsection{Keeping the time--frequency pattern of the record}\label{sec:pattern}

Instead of a random field with the statistics of the record, the field can follow the record itself
(\code{modulation = "pattern"}): the bursts then occur in the same bands and in the same order as in
the reference. The logarithm of the band power of the record, $\beta_j(t)$, is computed at every node
with the frame length that the synthesis uses at that node (so the time resolution is the one the
synthesis can reproduce), smoothed lightly, interpolated on the frame grid of the test and centred,
$\tilde\beta_j = \beta_j - \tavg{\beta_j}$. The field is
\begin{equation}
  m_j(t) = \frac{\exp\bigl(\gamma_j\tilde\beta_j(t)\bigr)}{\tavg{\exp(\gamma_j\tilde\beta_j)}} ,
\end{equation}
where the exponent $\gamma_j\ge0$ sets the depth: $\gamma_j=1$ reproduces the depth of the record, and
the design is free to change it to meet $\kappa_t$ and $D_t$. The map from $\gamma_j$ to the excess
$e_j = 3\Var(m_j)$ is tabulated on a grid of exponents and made monotone, so that the design can
work with $e_j$ as for the random field.

A test shorter than the record needs a \emph{mapping} of the reference time onto the test time:
\begin{itemize}[nosep]
\item \code{blocks} (default): blocks of a few seconds taken evenly across the whole record and
concatenated; every part of the record is represented and the burst durations are kept;
\item \code{compress}: the whole record squeezed into the test (the band power is averaged over each
test frame); the order of all events is kept and the bursts are shortened;
\item a test as long as the record uses it frame by frame (\code{identity}); a longer one repeats it
(\code{loop}).
\end{itemize}
The agreement between the band-power pattern of the result and the mapped reference is reported as
the median, over the nodes, of their correlation coefficient (1: same pattern).

\section{The MMVSD algorithm}\label{sec:mmvsdalg}

\subsection{Frames of several lengths}\label{sec:frames}

The modulated signal is built from short pieces of noise, \emph{frames}, each with its own spectrum
$m(f,t_\ell)$. A frame of length $L$ resolves frequency to $1/L$ and time to about $L$: long frames
are needed at low frequencies (an oscillator needs several cycles in a frame), short ones to follow
short bursts. The program therefore uses frames whose length grows towards low frequencies.

\emph{Requirements.} Node $i$ needs a hop (the spacing of consecutive frames, half the frame length)
of at most a third of its burst duration, and frames of at least 8 cycles, i.e.\ a hop of at least
$4/f_{\mathrm{n},i}$. The hop it gets is the larger of the two, $h_i = \max(T_{c,i}/3,\
4/f_{\mathrm{n},i})$, made non-increasing towards high frequencies (a cumulative minimum from the
lowest node up).

\emph{Groups.} The finest hop $h_0 = \min_ih_i$ (an even number of samples, at least 16) defines the
fine time grid. Every node is assigned to a group $g\in\{0,1,\dots\}$ with the hop $h_02^{g}$: the
longest power-of-two multiple of $h_0$ that does not exceed its need $h_i$, but long enough for the
cycle condition. At most six groups are used, the coarsest with at least 16 frames in the test, and the
test length is rounded to a whole number of coarse frames.

\emph{Crossovers.} Each group synthesises the nodes of its frequency range. Two adjacent groups are
joined at a boundary frequency $f_b$ (geometric mean of their edge nodes) with power-complementary
weights over 1/6 octave: with
\begin{equation*}
  t = \min\Bigl(1,\ \max\Bigl(0,\ 6\log_2\frac{f}{f_b} + \frac12\Bigr)\Bigr)
\end{equation*}
the lower group receives $1-\sin^2(\pi t/2)$ and the upper one $\sin^2(\pi t/2)$ of the power. Applied
recursively from the lowest boundary up, this defines weights $X_g(f)\ge0$ with $\sum_gX_g(f) = 1$
at every frequency.

\subsection{Synthesis of one group}

Let the group have the hop $h$ (in samples) and the frame length $L = 2h$. Frame $\ell$ covers the
samples $(\ell-1)h,\dots,(\ell+1)h-1$ (circularly) and is multiplied by the \emph{square-root Hann}
window $w(n) = \sin\bigl(\pi(n+\tfrac12)/L\bigr)$, $n = 0,\dots,L-1$. At any sample two frames overlap,
at positions $n$ and $n+h$ of their windows, and
\begin{equation}
  w(n)^2 + w(n+h)^2 = \sin^2\!\Bigl(\frac{\pi(n+\frac12)}{L}\Bigr) +
  \sin^2\!\Bigl(\frac{\pi(n+\frac12)}{L}+\frac{\pi}{2}\Bigr)
  = \sin^2(\cdot)+\cos^2(\cdot) = 1,
\end{equation}
because $\pi h/L = \pi/2$. The squared windows of overlapping frames therefore add up to one.

Each frame is a periodic noise built from its discrete Fourier coefficients on the lines $f_q =
q\,f_s/L$, $q = 0,\dots,h$:
\begin{equation}
  C_{\ell q} = \sqrt{\tfrac{Lf_s}{2}}\;\sqrt{X_g(f_q)}\;\sqrt{m_\ell(f_q)}\;c_{\ell q},
\end{equation}
where $c_{\ell q}$ are random coefficients with $\E|c_{\ell q}|^2=1$ (Section~\ref{sec:framenoise}) and
$m_\ell(f_q)$ is the field of the group (the fine field averaged over the $2^g$ fine frames of the
frame), interpolated linearly in $\log_2 f$ between the nodes. The inverse transform of length $L$
turns the coefficient $C$ of line $q$ into a cosine of amplitude $2|C|/L$, i.e.\ of power
$\tfrac12(2|C|/L)^2 = 2|C|^2/L^2$. With the factor $\sqrt{Lf_s/2}$ this power is $(f_s/L)\,X_gm_\ell|c|^2$,
and dividing by the line spacing $f_s/L$ gives the PSD $X_g(f)\,m_\ell(f)$ (on average over $c$):
with $m=1$ and $X_g=1$ the frame is white noise of unit one-sided PSD.

The frames of a group are independent and have zero mean, so at a given time the variances of the
two overlapping (windowed) frames add. Summing over the groups (also independent),
\begin{result}[Local spectrum of the synthesised noise]
\begin{equation}\label{eq:meff}
  m_{\mathrm{eff}}(f,t) = \sum_g X_g(f)\sum_\ell w_{g\ell}(t)^2\,m_{g\ell}(f),
\end{equation}
\end{result}
which is 1 for $m=1$ (since $\sum_\ell w^2 = 1$ and $\sum_g X_g = 1$) and has time average 1 when every
row of $m$ has mean 1. The model evaluates~\eqref{eq:meff} on the fine grid with two points per fine
hop ($\delta = h_0\Delta t/2$): between two frame centres the weights are $\cos^2(\pi u/2)$ and
$\sin^2(\pi u/2)$, $u\in(0,1)$ being the position between them, exactly the squared windows.

\emph{Shaping.} The sum of the groups, $v(t)$, is a noise whose local spectrum is
$m_{\mathrm{eff}}(f,t)$. Its discrete Fourier transform over the whole test, multiplied by
$\sqrt{P(f)} = \sqrt{\Pi(f)\,p(f)}$ on every line, gives the test signal: a zero-phase filter that
turns the unit spectrum into $P(f)$. The local spectrum of the result is $P(f)\,m_{\mathrm{eff}}(f,t)$,
which is the model~\eqref{eq:evol} with $m = m_{\mathrm{eff}}$.

\subsection{The frame noise}\label{sec:framenoise}

Two kinds of coefficients $c_{\ell q}$ are available.
\begin{itemize}
\item \emph{Gaussian frames}: complex normal coefficients (independent real and imaginary parts of
variance $\tfrac12$). Every frame is then a Gaussian noise, and so is the sum of frames and its
filtered version at every instant, for any frame length: the signal is conditionally Gaussian, as
the model assumes (Section~\ref{sec:multisine}).
\item \emph{Multisine frames}: coefficients of modulus 1 with random phases. Each frame has exactly
the band power of its spectrum, so the realisation scatters less around the model; but when a frame
contains only a few lines in the band of an oscillator ($L\,\fn/Q$ small, typically at low frequency
with short frames), the response within the frame is a sum of a few sinusoids whose amplitudes are
fixed: its peaks are lighter-tailed than Gaussian, and its damage is on average 20 to 40\,\% lower than the
Gaussian value at the same variance. The closed loop compensates with a higher level and a deeper
modulation.
\end{itemize}
The program uses Gaussian frames (\code{frame\_noise = "gaussian"}), so that the realisations
follow the statistics that the model assumes. On measured transport records (not distributed with
the package) they gave lower maximum FDS errors and lower $\kappa_z$ errors after the closed loop than
multisine frames; on synthetic records the two were within the scatter between seeds. The scatter
itself is a property of the target: with $b = 8$ the damage of a strongly modulated response rests on
its few largest bursts, and the damage of a single oscillator typically varies by a factor of about
1.5 between two noises of the same field (standard deviation of about 0.5 in $\ln D$ for a field of
depth $e = 3$ in a 60\,s test). This is the main reason for the closed loop of
Section~\ref{sec:mmvsdloop}.

\subsection{Design}\label{sec:design}

The design finds the nodal PSD factors $p_j$ and modulation depths $e_j$ for which the model meets the
targets, before any signal is synthesised. The unknowns are taken in logarithms, $y_j = \ln p_j$ and
$\chi_j = \ln e_j$, which keeps them positive and makes relative changes uniform. With the field
fixed, the model gives the residuals
\begin{equation}
  r_{D,i}(y) = \ln\frac{D_i(y)}{D_{t,i}}, \qquad r_{\kappa,i}(\chi) = \kappa_i(\chi) - \kappa_{t,i}.
\end{equation}
Since the kurtosis depends mainly on the depth and the damage on both, the two are solved
alternately, three times: first $\chi$ for $r_\kappa = 0$ with $y$ fixed, then $y$ for $r_D=0$ with
$\chi$ fixed.

Each problem is a regularised least-squares problem solved by the Levenberg--Marquardt method. For the
level, for example,
\begin{equation}
  \min_y\ \|r_D(y)\|^2 + \mu\,\|L(y - y_0)\|^2 ,
\end{equation}
where $L$ is the second-difference matrix, $(Ly)_i = y_i - 2y_{i+1} + y_{i+2}$, $\mu$ the weight of
the curvature penalty ($0.05$ for the level, $0.5$ for the depth) and $y_0$ the preferred solution
($y_0=0$ with the shape prior: no correction of the reference shape unless the targets require it).
Linearising $r_D(y+\delta) \approx r_D + J\delta$, with the Jacobian $J_{ij} = \partial r_{D,i}/\partial
y_j$, and setting the gradient of the objective to zero gives the step
\begin{equation}
  \bigl(J^{\top}J + \mu L^{\top}L + \lambda\,\mathrm{diag}(J^{\top}J+\mu L^{\top}L)\bigr)\,\delta
  = -\bigl(J^{\top}r_D + \mu L^{\top}L(y-y_0)\bigr).
\end{equation}
The damping $\lambda$ is divided by 3 after a step that lowers the objective and multiplied by 4 after
one that does not (then the step is retried). The depths are limited to $10^{-3}\le e_j\le e_{\max}$.
The level design starts from the median level of the stationary Gaussian solution of
Section~\ref{sec:inversion}, the depth design from $e_j = \kappa_{t,j}-3$.

\subsection{Exact Jacobians}\label{sec:jacobians}

The model~\eqref{eq:S}--\eqref{eq:Dmodel} is a composition of linear operators and pointwise
functions, so its derivatives can be computed exactly and cheaply.

\emph{Damage with respect to the level.} On the model grid of $N_t$ points, $D_i =
\frac{K^b}{C}\frac{T}{N_t}\sum_l d\bigl(S_i(t_l)\bigr)$. Let $g_{k,i}(t) = \partial d/\partial
S_k$ evaluated at $S_i(t)$. By~\eqref{eq:S}, $\partial S_{k,i}(t)/\partial y_j = W^{(k)}_{ij}
F_i[Y_j](t)$ with $Y_j(t) = p_jm_j(t)$ (because $\partial p_j/\partial y_j = p_j$). By the chain rule,
\begin{equation}
  \frac{\partial D_i}{\partial y_j} = \frac{K^b}{C}\frac{T}{N_t}\sum_k W^{(k)}_{ij}\sum_l g_{k,i}(t_l)\,F_i[Y_j](t_l).
\end{equation}
The inner sum is a scalar product between $g_{k,i}$ and a filtered signal. For any linear filter
$F$, $\sum_l g(t_l)\,F[Y](t_l) = \sum_l F^{\top}[g](t_l)\,Y(t_l)$, where the \emph{adjoint} $F^{\top}$ is
the same filter run backwards in time (in the Fourier domain: multiplication by the complex conjugate
of its response). Hence
\begin{result}[Damage Jacobian]
\begin{equation}
  \frac{\partial\ln D_i}{\partial y_j} = \frac{1}{D_i}\frac{K^b}{C}\frac{T}{N_t}
  \sum_k W^{(k)}_{ij}\sum_l F_i^{\top}[g_{k,i}](t_l)\;p_j\,m_j(t_l).
\end{equation}
\end{result}
The $4n$ adjoint filterings and one matrix product with the field give all $n^2$ entries at the cost
of a few model evaluations. The derivatives $g_{k,i}$ of the closed-form rate $d$ are computed by
one-sided differences with a relative step of $10^{-5}$.

\emph{Kurtosis with respect to the depth.} With $A_i = \tavg{s_i^2}$ and $B_i = \tavg{s_i}$,
$\kappa_i = 3A_i/B_i^2$. A change $\delta m_j(t)$ of the field row $j$ changes $s_i$ by $\delta s_i
= F_i[W^{(0)}_{ij}p_j\,\delta m_j]$. The mean $B_i$ does not change, because the field keeps
$\tavg{m_j} = 1$ and therefore $\tavg{\delta m_j} = 0$. The mean square changes by
$\delta A_i = 2\tavg{s_i\,\delta s_i} = 2W^{(0)}_{ij}p_j\tavg{F_i^{\top}[s_i]\,\delta m_j}$, so
\begin{equation}
  \frac{\partial\kappa_i}{\partial\chi_j} = \frac{6\,W^{(0)}_{ij}\,p_j}{B_i^2}\,
  \Bigl\langle F_i^{\top}[s_i]\;\frac{\partial m_j}{\partial\chi_j}\Bigr\rangle ,
\end{equation}
where $\partial m_j/\partial\chi_j$ is obtained by rebuilding the field with $e_j$ multiplied by
$e^{0.05}$ (the field depends smoothly on $e_j$). The same steps give the derivative of $\ln D_i$ with
respect to the depth, used by the closed loop.

\emph{Kurtosis with respect to the level.} Here $B_i$ changes too: $\partial B_i/\partial y_j =
W^{(0)}_{ij}p_j\tavg{m_j}$ and $\partial A_i/\partial y_j = 2W^{(0)}_{ij}p_j\tavg{F_i^\top[s_i]\,m_j}$,
and $\partial\kappa_i = 3\bigl(\partial A_i/B_i^2 - 2A_i\,\partial B_i/B_i^3\bigr)$.

\subsection{The closed loop on the realisation}\label{sec:mmvsdloop}

The design is synthesised once with frozen frame coefficients; from then on the signal is a smooth
function of $(p,e)$. Every iteration measures, on the realisation, the rainflow FDS $D_{m}$ and the
kurtosis of the responses $\kappa_{m}$ (both in one pass over the oscillators) and forms
\begin{equation}
  r_{D,i} = \ln\frac{D_{m,i}}{D_{t,i}},\qquad r_{\kappa,i} = w_\kappa\,\bigl(\bar\kappa_{m,i}-\kappa_{t,i}\bigr),
\end{equation}
with $w_\kappa = 0.15$ and $\bar\kappa_m$ the measured kurtosis averaged over three adjacent nodes.
The quality of a realisation is the score $\tfrac12\max_i|r_{D,i}| + \tfrac12q_{90}(|r_D|) +
\mathrm{mean}_i|r_{\kappa,i}|$, and the best realisation found is kept. Four kinds of steps are used.

\emph{Early re-draws.} The first realisation of the design can miss the target damage of one
oscillator by a large factor (several hundred per cent) while the median error is a few tens of per
cent: with $b$ between 5 and 10, one unusually large burst of the noise within the band of that
oscillator carries most of its damage. A correction of the level would then dig a notch into $P$ at
that node to compensate one rare event of the noise. Before any correction, therefore, while the worst
oscillator misses its damage by more than 50\,\%, the response $z$ of that oscillator is computed on
the realisation, the blocks of its group where $\sum|z|^b$ is largest are found, and the frames that
build them (the two largest blocks and their neighbours) receive new random coefficients on the lines
within $\pm1/4$ octave of the oscillator. The new draw is kept only if it lowers both the score and
the error of that oscillator; otherwise the old coefficients are restored. At most six draws are made,
and the phase ends after two useless draws in a row, because the remaining error is then a property
of the design and not of the noise.

\emph{Coupled Gauss--Newton steps} (first phase of the corrections). The model Jacobians carry the couplings between
oscillators (shared bands, the modulation of a neighbouring band, the envelope law), the measured
residuals what the model neglects. First the depth moves to correct the kurtosis,
\begin{equation}
  \delta\chi = -\bigl(J_{\kappa\chi}^{\top}J_{\kappa\chi} + \mu'L^{\top}L + \lambda\,\mathrm{diag}\bigr)^{-1}
  \bigl(J_{\kappa\chi}^{\top}r_\kappa + \mu'L^{\top}L\chi\bigr),
\end{equation}
then the level moves to correct the damage, including the change that the model predicts for the new
depth, $r_D + J_{D\chi}\delta\chi$:
\begin{equation}
  \delta y = -\bigl(J_{Dy}^{\top}J_{Dy} + \mu''L^{\top}L + \lambda\,\mathrm{diag}\bigr)^{-1}
  \bigl(J_{Dy}^{\top}(r_D + J_{D\chi}\delta\chi) + \mu''L^{\top}L(y-y_0)\bigr),
\end{equation}
with weak curvature weights ($\mu'' = 10^{-4}$, $\mu' = 0.05$) so that node-wise corrections remain
possible, and steps limited to $\pm0.5$ (depth) and $\pm0.3$ (level) in logarithm. After two steps
that do not improve the score, the loop switches to polishing.

\emph{Polishing.} Only the level moves, by the damped least-squares step
\begin{equation}
  \delta y = -w\,\bigl(J_{Dy}^{\top}J_{Dy} + \mu_p\,\mathrm{diag}(J_{Dy}^{\top}J_{Dy})\bigr)^{-1}J_{Dy}^{\top}r_D ,
  \qquad \mu_p = 0.5 .
\end{equation}
If the Jacobian were diagonal with the entries $b/2$ (each oscillator responding only to its own
node, $D\propto p^{b/2}$), this would be $\delta y_i = -w\,r_{D,i}/\bigl((b/2)(1+\mu_p)\bigr)$, i.e.\ the
node-wise step~\eqref{eq:gain} of version~2.0 ($\ln g_i^2 = (2w/b)\ln(D_t/D)$) reduced by $1+\mu_p$.
With the full Jacobian, the correction of an oscillator whose damage comes from a neighbouring band
goes to that band instead of driving its own node up or down. The relaxation $w$ adapts as in
Section~\ref{sec:closedloop}.

\emph{Re-draws.} The largest residual errors are usually a property of the realisation: the damage of
an oscillator rests on a few large cycles, and the level cannot fix it without disturbing its
neighbours. When the polishing stalls, the frames that carry most of the model damage of the worst
oscillator (the two largest and their neighbours, in its group) receive new random coefficients on
the lines within $\pm1/4$ octave of it, and a short branch of four polishing steps follows. The branch is
kept only if it beats the best realisation so far; otherwise the old coefficients are restored. At most
eight re-draws are tried.

\emph{Stopping.} The loop stops when the maximum FDS error is below the tolerance and the median
kurtosis error below 0.3, when the iteration budget is used, or when the score has not improved by
1\,\% for ten iterations.

\emph{A second strategy.} The coupled loop corrects kurtosis and damage together and usually
converges in a few tens of iterations. Where the damage of a few oscillators rests on a few bursts it
can stall at a compromise between neighbouring oscillators, because the model Jacobian then
describes the realisation poorly. On request (option \code{fallback}), when its maximum FDS error ends
above twice the tolerance, the program runs a second closed loop from the same design and the same
noise with node-wise steps only (the step of Section~\ref{sec:closedloop}, the depth kept at the
design), which progresses more slowly but steadily, and keeps the realisation with the better score.
The option is off by default: on measured transport records the coupled loop was better in every
case, and the second loop only doubled the computing time. The program warns when the final level differs from the design by more than
a factor of 8 over a band (a systematic drift, typically at oscillators whose damage is produced by
another band), and when the target kurtosis could not be reached with the maximum depth.

\section{Sine-on-random blocks}\label{sec:lines}

\subsection{The problem}

Rotating machines (engines, pumps, fans) add tonal lines to the random vibration, at the rotation
frequency and its orders. In a record these lines are often present only in some parts (idle, cruise)
and their frequency drifts with the speed. A line is not a narrow random resonance: its phase is
coherent over many periods, its amplitude is nearly constant, and a sinusoidal response is far less
damaging than a Gaussian response of the same variance (a sine has no rare large cycles: its damage
ratio to a narrow-band Gaussian response of equal variance is $1/\Gamma(1+b/2)$, which is $1/24$ for
$b=8$, from~\eqref{eq:rayleighmoment} with $A = \sqrt2\sigma$ fixed). Treating the lines as part of a
random process therefore misrepresents the damage. MMVSD detects the blocks of the record where lines
stand on the random background (\emph{sine-on-random blocks}), reproduces the lines as deterministic
sines, and designs the random part for the rest.

\subsection{Detection block by block}

The record is cut into blocks of 5\,s. In every block the PSD is estimated with a resolution of
0.5\,Hz (Hann windows, 75\,\% overlap) and the line detector of Section~\ref{sec:sor2} proposes
candidate peaks. A candidate is accepted only if it passes a \emph{phase coherence} test.

\emph{Coherence.} The frequency of the candidate is followed frame by frame inside the block (the
maximum of each short-time spectrum near the candidate, refined by fitting a parabola to the logarithm
of the three largest values), and a straight line $f(t)$ is fitted to these values. The block is then
demodulated along this frequency, $y(t) = x(t)\,e^{-2\pi\ii\int f\,\dd t}$, and low-pass filtered to
$\pm1$\,Hz, which leaves the \emph{complex envelope} $c(t)$ of whatever lies near the candidate. For a
sine $A\cos(\varphi(t)+\psi)$ with $\varphi = 2\pi\int f\dd t$,
\begin{equation}
  x\,e^{-\ii\varphi} = \tfrac{A}{2}e^{\ii\psi} + \tfrac{A}{2}e^{-\ii\psi}e^{-2\ii\varphi},
\end{equation}
and the low-pass filter removes the second term (at twice the line frequency): $c = \tfrac
A2e^{\ii\psi}$ is constant. For narrow-band noise $c(t)$ is a random phasor. With $N$ samples of $c$,
the coherence is
\begin{equation}
  \Phi = \max_{|\delta|\le r/4}\frac{\bigl|\sum_l c_l\,e^{-2\pi\ii\delta t_l}\bigr|^2}{N\sum_l|c_l|^2}\in[0,1],
\end{equation}
where the small residual frequency offset $\delta$ (up to a quarter of the resolution $r$) absorbs
the error of the frequency fit. By the Cauchy--Schwarz inequality $|\sum_lc_l|^2\le N\sum_l|c_l|^2$, with
equality when $c$ is constant: a sine gives $\Phi\approx1$ (slightly less because of the background in
the filter band). Noise of bandwidth $B$ observed for the block duration $T_b$ has a complex envelope
that stays correlated for about $1/B$; the $N$ samples form about $N_{\mathrm{ind}} = BT_b$
independent groups of $N/N_{\mathrm{ind}}$ nearly equal values. Each group sums to about
$(N/N_{\mathrm{ind}})c_g$, the groups add with random phases, so $\E|\sum_lc_l|^2\approx
N_{\mathrm{ind}}(N/N_{\mathrm{ind}})^2\E|c|^2 = N^2\E|c|^2/N_{\mathrm{ind}}$, while $N\sum|c|^2\approx
N^2\E|c|^2$: noise gives $\Phi\approx1/(BT_b)$. A random resonance with $Q = 60$ at 70\,Hz ($B\approx
1.2$\,Hz) over 5\,s gives $\Phi\approx0.17$; a line gives 0.9 or more.

\emph{Tracks.} A candidate with $\Phi\ge0.75$ is a line; candidates of consecutive blocks within 5\,\%
in frequency are linked into \emph{tracks}. A track is kept if at least two of its blocks are strongly
coherent ($\Phi\ge0.75$); between them it may pass through weakly coherent blocks ($\Phi\ge0.5$), so
that a line with a wandering speed or a low signal-to-noise ratio is not cut (hysteresis). A random
resonance that passes the test by chance in one block is dropped.

\subsection{Extraction}

Each track is followed frame by frame in its active intervals (consecutive detected blocks, gaps of up
to two blocks bridged), its frequency is median-filtered over three frames against spurious peaks,
and the record is demodulated along the track's own phase and low-pass filtered (to half the
resolution). By the computation above, the low-pass result is the complex envelope $c(t)$ of the line,
and the line itself is
\begin{equation}
  \ell(t) = 2\,\mathrm{Re}\bigl(c(t)\,e^{\ii\varphi(t)}\bigr),
\end{equation}
tapered at the edges of the intervals. The sum of the lines is subtracted from the record; the
\emph{residual} is the random part of the record. The track keeps its frequency $f(t)$ and amplitude
$a(t) = 2|c(t)|$ on a time grid of the record.

\subsection{The lines of the test}\label{sec:testlines}

The reference time is mapped onto the test time with the mapping of the time--frequency pattern
(Section~\ref{sec:pattern}); with the random field, \code{compress} is used, so that every tonal block
of the record appears in the test, in its order. The amplitude and frequency of each track are read
along the mapped time, smoothed over 0.1\,s, and the track is played as a sine with a random initial
phase.

\emph{Gain of a track.} The damage that a line with a slowly varying amplitude $a(t)$ produces in an
oscillator is the sum of the damages of its cycles, $\frac{K^b}{C}\int f\,(|H|a)^b\,\dd t$. If the
test plays the same frequency path, the factors $f|H|^b$ are the same in the record and in the test,
and the life damage is kept when the test amplitude $g_ja_j$ satisfies
$\int_{\mathrm{test}}(g_ja_j)^b\dd t = \frac{\Tl}{T_{\mathrm{ref}}}\int_{\mathrm{ref}}a_j^b\dd t$:
\begin{result}[Gain of a track]
\begin{equation}
  g_j = \left(\frac{\Tl}{T_{\mathrm{ref}}}\,\frac{\int_{\mathrm{ref}}a_j(t)^b\,\dd t}{\int_{\mathrm{test}}a_j(t)^b\,\dd t}\right)^{1/b}.
\end{equation}
\end{result}
For a stationary line, $\int_{\mathrm{ref}}a^b = A^bT_{\mathrm{ref}}$ and $\int_{\mathrm{test}}a^b = A^bT$,
and $g = (\Tl/T)^{1/b}$, the sine-on-random rule~\eqref{eq:sorgain}. A track whose blocks fall outside
the mapped part of the record is reported; its damage is then carried by the random part.

\subsection{The random part}

Line damage and random damage do not add: a line rides on the random peaks and changes their ranges.
Rather than assuming additivity, MMVSD designs the random part for the \emph{measured} FDS and response
kurtosis of the residual, scaled to the life, $D_{\mathrm{res}}\Tl/T_{\mathrm{ref}}$ and
$\max(\kappa_{\mathrm{res}},3)$, with the PSD of the residual as shape prior and, in pattern mode, the
band-power pattern of the residual. The closed loop then measures the complete test (random part plus
lines) against the complete target, with the lines kept fixed.

\emph{Kurtosis of a response with lines.} Given the field, the response of an oscillator is a
Gaussian part $g$ (variance $s$) plus the responses of the lines, $\sum_j u_j^{1/2}\sqrt2\cos\theta_j$,
where $u_j$ is the response variance of line $j$ and the phases $\theta_j$ are independent and uniformly
distributed. For one line of amplitude $\alpha = \sqrt{2u}$: $\E[\alpha^4\cos^4\theta] =
\alpha^4\cdot\tfrac38 = \tfrac38(2u)^2 = \tfrac32u^2$, using $\cos^4\theta = \tfrac38 +
\tfrac12\cos2\theta + \tfrac18\cos4\theta$. For the sum $X = \sum_jX_j$ of independent zero-mean lines,
the fourth moment keeps the terms $X_j^4$ and the pairs $X_j^2X_l^2$ ($\binom42 = 6$ orderings,
counted once for each unordered pair, i.e.\ $3$ per ordered pair $j\neq l$):
\begin{equation}
  \E[X^4] = \sum_j\tfrac32u_j^2 + 3\sum_{j\ne l}u_ju_l = 3U^2 - \tfrac32\sum_ju_j^2,\qquad U = \sum_ju_j ,
\end{equation}
because $3\sum_{j\ne l}u_ju_l = 3(U^2 - \sum_ju_j^2)$. Adding the independent Gaussian part, whose odd
moments vanish, $\E[(g+X)^4] = \E g^4 + 6\E g^2\E X^2 + \E X^4 = 3s^2 + 6sU + 3U^2 -
\tfrac32\sum u_j^2$:
\begin{result}[Fourth moment of a Gaussian response with lines]
\begin{equation}
  \E[z^4\mid m] = 3(s+U)^2 - \tfrac32\sum_ju_j^2 ,
  \qquad \kappa = \frac{3\tavg{(s+U)^2} - \tfrac32\tavg{\sum_ju_j^2}}{\tavg{s+U}^2}.
\end{equation}
\end{result}
The line variances $u_j(t) = \tfrac12a_j(t)^2|H_d(f_j(t))|^2|H_{\mathrm{pre}}|^4$ are filtered by the
envelope law like $s$. The kurtosis model of the closed loop uses this formula; the damage of the
lines has no reliable closed form on a broad-band random response, so it is not modelled but measured.

\emph{Damage elasticity.} Without lines, multiplying the PSD by $\lambda$ multiplies the amplitudes by
$\sqrt\lambda$ and the damage by $\lambda^{b/2}$: $\partial\ln D/\partial\ln p = b/2$. With lines part of
the damage is fixed and the elasticity $\eta_i$ is lower. It is measured on the first realisation by
amplifying the random part by $1.1$ (its PSD by $1.1^2$):
\begin{equation}
  \eta_i = \frac{\ln\bigl(D_i(1.1\,x_r + x_\ell)/D_i(x_r+x_\ell)\bigr)}{2\ln1.1},
\end{equation}
limited to $[0.1\,b/2,\ b/2]$, and the damage rows of the closed-loop Jacobians are multiplied by
$\eta_i/(b/2)$. Where the lines alone exceed the target damage, the program reports that those
frequencies are over-tested, as the sine-on-random method of version~2.0 does.

\section{From the theory to the settings}\label{sec:practice}

\subsection{The damage model}

The damage model (natural frequency grid, $Q$, $b$, $C$, $K$, pre-filter) is shared by all
computations of a workspace. The FDS of a record is computed by default with the exact procedure of
Section~\ref{sec:fds} (\code{time}); \code{frequency} uses the Gaussian formula~\eqref{eq:dgauss} on the
Welch PSD, \code{fast} a self-calibrated approximation. If $f_{\max}>f_s/2.5$ the record is first
up-sampled by $5/4$, so that every oscillator has at least 2.5 samples per period. The exact
computation also returns the variance and the kurtosis of every response, $\kappa_z(\fn)$, in the same
pass; the program computes the oscillators in parallel (worker processes for long records), and the
results do not depend on the number of workers.

\subsection{Choosing a method}

\begin{center}
\begin{tabular}{p{0.52\textwidth}p{0.38\textwidth}}
\toprule
Situation & Method \\
\midrule
Only the damage matters; a Gaussian test is acceptable & \code{gaussian} \\
Leptokurtic record, one kurtosis for the whole signal & \code{kfds} \\
The test must keep the RMS of the field & \code{rf} \\
The test must keep the PSD (shape) of the field & \code{pf}, \code{psf} \\
Stationary tonal lines over the whole record & \code{sor} \\
Kurtosis of the responses different in different bands; lines in some blocks only; bursts at the
right time & \code{mmvsd} \\
Peak severity (shock-like environments) & \code{mrs} \\
Kurtosis control without damage equivalence & \code{mlv}, \code{vsd}, \code{pm} \\
\bottomrule
\end{tabular}
\end{center}

\subsection{MMVSD settings}

\begin{center}
\begin{tabular}{p{0.24\textwidth}p{0.20\textwidth}p{0.46\textwidth}}
\toprule
Interface & Command line & Meaning \\
\midrule
Modulation & \code{--modulation} & \code{random}: the field of Section~\ref{sec:field} with the burst
statistics of the target; \code{pattern}: the time--frequency pattern of the record
(Section~\ref{sec:pattern}) \\
Shorter test: mapping & \code{--mapping} & \code{blocks} or \code{compress} (Section~\ref{sec:pattern}) \\
Block length & \code{--block} & length of the blocks of the \code{blocks} mapping (default: 4\,s or five
burst durations) \\
Sine-on-random blocks & \code{--lines} & \code{lines}: detected lines as deterministic sines
(Section~\ref{sec:lines}); \code{random}: detected but left in the random part; \code{off} \\
Burst width / duration & \code{--freq-corr}, \code{--time-corr} & $B$ in octaves and $T_c$ in
seconds (Section~\ref{sec:structure}); estimated from the record by default \\
Max.\ modulation excess & \code{--max-excess} & cap $e_{\max}$ of the depth (default 30) \\
Design anchor & \code{--anchor} & pull of the level towards the design at the oscillators whose
damage comes from another band (default 0) \\
Max.\ corrections, tolerance & \code{--iterations}, \code{--tol} & budget and stopping tolerance of
the closed loop \\
Seed & \code{--seed} & the random numbers of the field and of the frames: the same seed gives the
same profile \\
\bottomrule
\end{tabular}
\end{center}

For a mission or a specification there is no single record: the target $\kappa_z$ is the mixture
of Section~\ref{sec:missions}, the burst statistics are estimated on the longest situation, and only
the random field is available. Lines are detected on single records only.

\subsection{Reading the results}

\begin{itemize}
\item \emph{FDS error}: the ratio of the rainflow FDS of the profile (times its repetitions) to the
target, node by node; the maximum and the median of $|{\rm ratio}-1|$ and the share of nodes within
10\,\% are reported. The maximum is governed by the few oscillators whose damage rests on a few large
cycles (low natural frequencies, strong modulation); the median shows the overall agreement.
\item \emph{$\kappa_z$ error}: the median of $|\kappa_{z,\mathrm{profile}}-\kappa_{z,\mathrm{target}}|$
over the nodes, computed for every method (it shows, for instance, that a Gaussian test has
$\kappa_z\approx3$ everywhere).
\item \emph{MRS ratio}: above 1 the test produces larger peak responses than the target; check it
against the strength of the component.
\item \emph{PSD deviation}: the RMS, over the band, of the ratio in dB between the PSD of the profile and
that of the target.
\item \emph{Pattern correlation} (pattern mode): Section~\ref{sec:pattern}.
\item \emph{Warnings}: a level drifted far from the design over a band, a kurtosis target that cannot
be reached with the maximum depth, bursts shorter than two response times (filtered by the oscillators,
Section~\ref{sec:papoulis}), lines that exceed the target damage alone or that are not in the mapped
part of the record.
\end{itemize}

\subsection{Limits}

The FDS and the response kurtosis describe the damage of linear, lightly damped oscillators with a
Basquin law; they are only as good as the chosen $Q$ and $b$ (the program computes the sensitivity of
a result to other values). The kurtosis of the responses is a fourth-moment description: two signals
with the same $\kappa_z$ can still differ in their largest peaks, which the MRS ratio shows. The
model of Section~\ref{sec:mmvsdmodel} assumes that the modulation varies slowly along frequency compared
with the oscillator bandwidth; very narrow bursts in frequency are smoothed. Finally, a test that is
much shorter than the life compresses the bursts: their duration is kept in the \code{blocks} mapping,
at the price of leaving out parts of the record.

\section{Diagnostics and checks}\label{sec:diagnostics}

The interface offers a set of diagnostics of a single signal and of a result. This section defines
them.

\subsection{Statistics over time and amplitude}

\emph{Running RMS and kurtosis.} The signal is cut into consecutive windows of $L_w$ seconds; in each
window the RMS and the kurtosis~\eqref{eq:kurtosis} are computed. A Gaussian stationary signal gives a
nearly constant RMS and a kurtosis near 3 in every window; bursts show up as RMS peaks. The
distribution of the 1-second RMS over the record (shown for a result and its target) describes the
burst structure that kFDS and MMVSD reproduce.

\emph{Probability density.} The histogram of the samples is compared with the Gaussian density with
the same mean and standard deviation, on a logarithmic scale where the tails are visible.

\subsection{Time--frequency diagnostics}

\emph{Spectrogram.} The short-time Fourier transform (STFT) cuts the signal into overlapping
windowed segments (Hann window, 50\,\% overlap) and computes the PSD of each, giving the power as a
function of time and frequency.

\emph{Spectral kurtosis.} Let $X(f,t)$ be the STFT coefficient at frequency $f$ of the segment at time
$t$. The spectral kurtosis is
\begin{equation}
  SK(f) = \frac{\tavg{|X(f,\cdot)|^4}}{\tavg{|X(f,\cdot)|^2}^2} - 2 .
\end{equation}
For a stationary Gaussian signal, $X = U + \ii V$ with $U$ and $V$ independent zero-mean Gaussian
variables of equal variance $\sigma^2/2$. Then $|X|^2 = U^2+V^2$ has the exponential distribution with
mean $\sigma^2$ (the squared length of a two-dimensional Gaussian vector; with
Section~\ref{sec:rayleigh}, $p(A) = (2A/\sigma^2)e^{-A^2/\sigma^2}$ for $A = |X|$, and the
substitution $u = A^2$ gives $p(u) = e^{-u/\sigma^2}/\sigma^2$), whose second moment is
$\int_0^\infty u^2e^{-u/\sigma^2}\dd u/\sigma^2 = 2\sigma^4$. Hence $SK = 2\sigma^4/\sigma^4 - 2 = 0$. A
band whose power changes from segment to segment (bursts) has $SK>0$; a stationary sine, whose
$|X|$ is constant, has $SK = 1 - 2 = -1$. The \emph{kurtogram} shows $SK$ as a function of frequency
and of the window length, since short bursts are seen only with short windows.

\emph{Response of one oscillator.} The pseudo-acceleration $\wn^2z(t)$ of a chosen oscillator, with
its statistics, shows directly how the input is filtered (Section~\ref{sec:papoulis}). The curves
of the kurtosis and standard deviation of all responses versus $\fn$ are the quantities $\kappa_z(\fn)$
and $\sigma_{\wn^2 z}(\fn)$ of Part~I, computed in one pass over the oscillators.

\subsection{Approximate FDS}

Besides the exact rainflow FDS and the Gaussian formula, the program offers the fast approximation of
the dissertation. For a narrow-band response, the large values of $z$ occur near the peaks of the
cycles, and a cycle of amplitude $A$ contributes to $\sum_kz_k^b\Delta t$ (over the samples where $z$ is
positive) an amount proportional to $A^b/\fn$. Hence
\begin{equation}
  D(\fn) \approx h(\fn)\,\frac{K^b}{C}\,\fn\,\Delta t\sum_{k:\ z_k>\rho\,z_{\max}} z_k^b ,
  \qquad \rho = 400^{-1/b},
\end{equation}
where only the samples above the fraction $\rho$ of the maximum are summed (smaller cycles contribute
little to the damage when $b$ is large) and $h(\fn)$ absorbs the proportionality constant and the
residual effects of the cycle shapes. The program calibrates $h(\fn)$ on the first 20\,s of the
signal itself against the exact rainflow damage.

\subsection{Sensitivity to $b$ and $Q$}

A profile is designed with one damage model. If the true exponent of the S--N curve or the damping
of the component differ, the profile and the target no longer have the same FDS. The program
recomputes the FDS of both with other values of $b$ or $Q$ and shows their ratio: a ratio that stays
close to 1 means that the equivalence does not depend on the exact damage model.

\subsection{Matching the amplitude distribution}

A profile can be given exactly the amplitude distribution of the target by a monotone mapping of its
samples: the sample of rank $r$ (out of $N$) is replaced by the quantile of the target values at the
probability $(r+\tfrac12)/N$. The order of the samples, hence the time structure, is kept; the PSD and
the FDS change slightly, so the program measures them again after the mapping.

\subsection{Shaker limits}

A shaker has limits of acceleration, velocity and stroke (peak-to-peak displacement). Velocity and
displacement are obtained from the acceleration in the frequency domain, dividing every line by
$\ii\omega$ (velocity) or by $-\omega^2$ (displacement), after a high-pass filter with the gain
$1/\sqrt{1+(f_{\mathrm{hp}}/f)^8}$ (the magnitude of a fourth-order Butterworth filter): without it
the lowest lines, divided by a very small $\omega^2$, would dominate the displacement. If the stroke
is exceeded, the program searches the lowest high-pass frequency that brings it within the limit;
applying that filter in the pre-processing changes the profile only below that frequency.


\section*{Acknowledgements}
\addcontentsline{toc}{section}{Acknowledgements}

The methods were developed during my PhD at the University of Bologna; I thank Prof.\ Marco Troncossi
and Easting s.r.l.s.\ for our joint work in those years. This manual and the Python implementation
were written with the help of Claude (Anthropic); all derivations were checked against the
implementation and numerically.

\appendix
\section{Lalanne's operator in closed form}\label{app:lalanne}

This appendix integrates~\eqref{eq:W}: the response variance of an oscillator to a PSD that is
linear between nodes.

\subsection*{A.1 Dimensionless form}

Fix the oscillator of natural frequency $\fn$ and write every frequency relative to it, $h = f/\fn$.
Then $\omega = h\wn$, and~\eqref{eq:H2} becomes
\begin{equation}
  |H|^2 = \frac{1}{(\wn^2-h^2\wn^2)^2 + (2\zeta\wn\cdot h\wn)^2} = \frac{1}{\wn^4\,\Delta(h)},
  \qquad \Delta(h) = (1-h^2)^2 + (2\zeta h)^2 .
\end{equation}
With $\dd f = \fn\,\dd h$, the variance is
\begin{equation}\label{eq:A1}
  \sigma_z^2 = \int|H|^2G\,\dd f = \frac{\fn}{\wn^4}\int\frac{G(\fn h)}{\Delta(h)}\,\dd h .
\end{equation}

\subsection*{A.2 Factorisation of $\Delta$}

Expanding, $\Delta(h) = 1 - 2h^2 + h^4 + 4\zeta^2h^2 = (h^2+1)^2 - (4-4\zeta^2)h^2$. With
\begin{equation}
  a = 2\sqrt{1-\zeta^2}\qquad (a^2 = 4-4\zeta^2),
\end{equation}
this is a difference of squares, $\Delta = (h^2+1)^2-(ah)^2$:
\begin{equation}
  \Delta(h) = (h^2+ah+1)(h^2-ah+1).
\end{equation}
Both factors are positive: $h^2\pm ah+1 = (h\pm a/2)^2 + 1 - a^2/4$ and $1-a^2/4 = 1-(1-\zeta^2) =
\zeta^2$, so
\begin{equation}\label{eq:A2}
  h^2\pm ah+1 = \Bigl(h\pm\frac a2\Bigr)^2 + \zeta^2 .
\end{equation}

\subsection*{A.3 Partial fractions}

Since $h\to-h$ exchanges the two factors, look for
\begin{equation}
  \frac{1}{\Delta} = \frac{\alpha h+\beta}{h^2+ah+1} + \frac{-\alpha h+\beta}{h^2-ah+1}.
\end{equation}
Over the common denominator the numerator is
$(\alpha h+\beta)(h^2-ah+1) + (-\alpha h+\beta)(h^2+ah+1)$. Expanding both products,
$\alpha h^3 - a\alpha h^2 + \alpha h + \beta h^2 - a\beta h + \beta$ and
$-\alpha h^3 - a\alpha h^2 - \alpha h + \beta h^2 + a\beta h + \beta$; their sum is
$(2\beta - 2a\alpha)h^2 + 2\beta$. It must equal 1 for all $h$: $2\beta = 1$ and $2\beta = 2a\alpha$,
so $\beta = \tfrac12$ and $\alpha = \tfrac{1}{2a}$:
\begin{equation}
  \frac{1}{\Delta} = \frac{\frac{h}{2a}+\frac12}{h^2+ah+1} + \frac{-\frac{h}{2a}+\frac12}{h^2-ah+1}.
\end{equation}

\subsection*{A.4 The two basic integrals}

\emph{First integral.} In the first fraction write the numerator as a multiple of the derivative of
the denominator plus a constant: $\frac{h}{2a}+\frac12 = \frac{2h+a}{4a} + \frac14$ (indeed
$\frac{2h}{4a}+\frac{a}{4a}+\frac14 = \frac{h}{2a}+\frac12$). The derivative of $h^2+ah+1$ is $2h+a$,
and by~\eqref{eq:A2}
\begin{equation}
  \int\frac{\dd h}{h^2+ah+1} = \int\frac{\dd h}{(h+\frac a2)^2+\zeta^2} = \frac1\zeta\arctan\frac{h+\frac a2}{\zeta}
  = \frac1\zeta\arctan\frac{2h+a}{2\zeta}.
\end{equation}
Hence $\int\frac{\frac{h}{2a}+\frac12}{h^2+ah+1}\dd h = \frac{1}{4a}\ln(h^2+ah+1) +
\frac{1}{4\zeta}\arctan\frac{2h+a}{2\zeta}$. The second fraction is treated in the same way with
$-\frac{h}{2a}+\frac12 = -\frac{2h-a}{4a}+\frac14$, giving $-\frac{1}{4a}\ln(h^2-ah+1) +
\frac{1}{4\zeta}\arctan\frac{2h-a}{2\zeta}$. Adding,
\begin{equation}\label{eq:A3}
  \int\frac{\dd h}{\Delta} = \frac{1}{4a}\ln\frac{h^2+ah+1}{h^2-ah+1}
  + \frac{1}{4\zeta}\left[\arctan\frac{2h+a}{2\zeta}+\arctan\frac{2h-a}{2\zeta}\right].
\end{equation}

\emph{Second integral.} Substitute $v = h^2$, $\dd v = 2h\,\dd h$:
\begin{equation}
  \int\frac{h\,\dd h}{\Delta} = \frac12\int\frac{\dd v}{(1-v)^2+4\zeta^2v} = \frac12\int\frac{\dd v}{v^2+(4\zeta^2-2)v+1}.
\end{equation}
Complete the square: $(v+2\zeta^2-1)^2 = v^2 + (4\zeta^2-2)v + (2\zeta^2-1)^2$, and
$1-(2\zeta^2-1)^2 = 1 - 4\zeta^4 + 4\zeta^2 - 1 = 4\zeta^2(1-\zeta^2) = \zeta^2a^2$, so the denominator
is $(v+2\zeta^2-1)^2 + \zeta^2a^2$ and
\begin{equation}\label{eq:A4}
  \int\frac{h\,\dd h}{\Delta} = \frac{1}{2\zeta a}\arctan\frac{h^2+2\zeta^2-1}{\zeta a}.
\end{equation}
Both antiderivatives are continuous in $h$. It is convenient to normalise them by $4\zeta/\pi$:
\begin{align}
  I_0(h) &= \frac{4\zeta}{\pi}\int\frac{\dd h}{\Delta}
  = \frac{\zeta}{\pi a}\ln\frac{h^2+ah+1}{h^2-ah+1} + \frac1\pi\left[\arctan\frac{2h+a}{2\zeta}+\arctan\frac{2h-a}{2\zeta}\right],\\
  I_1(h) &= \frac{4\zeta}{\pi}\int\frac{h\,\dd h}{\Delta} = \frac{2}{\pi a}\arctan\frac{h^2+2\zeta^2-1}{\zeta a}.
\end{align}
(The program writes $I_1$ as $\frac{2}{\pi a}[\arctan\frac{2h-a}{2\zeta}-\arctan\frac{2h+a}{2\zeta}]$,
which has the same derivative and therefore differs only by a constant, irrelevant in definite
integrals.)

\emph{White noise.} At $h=0$: $\ln1 = 0$ and $\arctan\frac{a}{2\zeta}+\arctan\frac{-a}{2\zeta} = 0$,
so $I_0(0) = 0$. For $h\to\infty$ the logarithm tends to $\ln1=0$ and each arctangent to $\pi/2$, so
$I_0(\infty) = 1$. Therefore $\int_0^\infty\dd h/\Delta = \pi/(4\zeta)$, and by~\eqref{eq:A1} with
$G = G_0$
\begin{equation}
  \int_0^\infty|H|^2\dd f = \frac{\fn}{\wn^4}\,\frac{\pi}{4\zeta} = \frac{\wn}{2\pi}\cdot\frac{\pi}{4\zeta\wn^4}
  = \frac{1}{8\zeta\wn^3},
\end{equation}
the result used in~\eqref{eq:miles}.

\subsection*{A.5 A PSD linear between nodes}

Let the nodes, in units of the oscillator's natural frequency, be $h_l = f_{\mathrm{n},l}/\fn$. On the
segment $[h_l, h_{l+1}]$ of width $\delta_l = h_{l+1}-h_l$ the PSD is
\begin{equation}
  G = G_l\,\frac{h_{l+1}-h}{\delta_l} + G_{l+1}\,\frac{h-h_l}{\delta_l}.
\end{equation}
With the increments $\Delta I_{0,l} = I_0(h_{l+1})-I_0(h_l)$ and $\Delta I_{1,l} = I_1(h_{l+1})-I_1(h_l)$,
\begin{equation}
  \frac{4\zeta}{\pi}\int_{h_l}^{h_{l+1}}\frac{G}{\Delta}\dd h
  = G_l\,\frac{h_{l+1}\Delta I_{0,l}-\Delta I_{1,l}}{\delta_l} + G_{l+1}\,\frac{\Delta I_{1,l}-h_l\Delta I_{0,l}}{\delta_l}.
\end{equation}
The node $j$ appears as the right end of segment $j-1$ and as the left end of segment $j$. Collecting
its two coefficients and multiplying by $\frac{\fn}{\wn^4}\frac{\pi}{4\zeta} = \frac{1}{8\zeta\wn^3}$
(from~\eqref{eq:A1}),
\begin{result}[Lalanne's operator]
\begin{equation}
  W_{ij} = \frac{1}{8\zeta\wn^3}\left[\frac{\Delta I_{1,j-1}-h_{j-1}\Delta I_{0,j-1}}{\delta_{j-1}}
  + \frac{h_{j+1}\Delta I_{0,j}-\Delta I_{1,j}}{\delta_j}\right],
\end{equation}
\end{result}
where $\wn$, $h$ and the increments refer to the oscillator $i$, and the first (second) term is absent
for the first (last) node. The program evaluates this for all pairs $(i,j)$ at once. The moment
operators of MMVSD (Section~\ref{sec:operators}) include the shape prior, the digital response and the
edge extension, which have no closed form; they are integrated numerically, and with $\Pi = 1$ and the
analogue response they reproduce this matrix to better than $10^{-3}$ (checked in the tests).

\section{The ramp-invariant oscillator filter}\label{app:smallwood}

This appendix derives the coefficients of the recursion~\eqref{eq:recursion} (the ramp-invariant
filter, after Smallwood). Throughout, $\omega = \wn$, $\wdm = \omega\sqrt{1-\zeta^2}$ is the damped
natural angular frequency, and the symbol $q$ denotes the variable of the $z$-transform (the letter $z$
is the relative displacement).

\subsection*{B.1 The input between samples}

Assume that the base acceleration varies linearly between consecutive samples. With the triangle
function $\Lambda(\theta) = \max(0,1-|\theta|)$, which is 1 at $\theta=0$, 0 at $\theta=\pm1$ and linear
in between,
\begin{equation}
  x(t) = \sum_k x_k\,\Lambda\!\left(\frac{t}{\Delta t}-k\right)
\end{equation}
is exactly the linear interpolation of the samples. The oscillator is linear and time invariant, so
its response at the sampling instants is
\begin{equation}\label{eq:B1}
  z_l = z(l\Delta t) = \sum_k x_k\,\gamma_{l-k},
\end{equation}
where $\gamma_l$ is the response, at $t = l\Delta t$, to the single triangle $\Lambda(t/\Delta t)$ centred
at $t=0$, starting from rest.

\subsection*{B.2 The triangle as a combination of ramps}

Let $r(t) = t$ for $t\ge0$ and $0$ for $t<0$ (a unit ramp starting at 0). Then
\begin{equation}
  \Lambda(t/\Delta t) = \frac{1}{\Delta t}\bigl[r(t+\Delta t) - 2r(t) + r(t-\Delta t)\bigr].
\end{equation}
Check on the four intervals: for $t\le-\Delta t$ all ramps are zero; for $-\Delta t\le t\le0$ only
the first is active and gives $(t+\Delta t)/\Delta t$, rising from 0 to 1; for $0\le t\le\Delta t$ it
gives $(t+\Delta t-2t)/\Delta t = (\Delta t-t)/\Delta t$, falling from 1 to 0; for $t\ge\Delta t$ it
gives $(t+\Delta t-2t+t-\Delta t)/\Delta t = 0$. By linearity, if $\rho(t)$ is the response to the ramp
$r(t)$,
\begin{equation}\label{eq:B2}
  \gamma_l = \frac{1}{\Delta t}\bigl[\rho_{l+1} - 2\rho_l + \rho_{l-1}\bigr],\qquad \rho_l = \rho(l\Delta t),
\end{equation}
with $\rho_l = 0$ for $l\le0$.

\subsection*{B.3 The ramp response}

$\rho$ solves $\ddot\rho + 2\zeta\omega\dot\rho + \omega^2\rho = t$ for $t\ge0$ with $\rho(0) =
\dot\rho(0) = 0$. A particular solution is $\rho_p = t/\omega^2 - 2\zeta/\omega^3$: indeed $\dot\rho_p =
1/\omega^2$, $\ddot\rho_p = 0$, and $2\zeta\omega/\omega^2 + \omega^2(t/\omega^2 - 2\zeta/\omega^3) =
2\zeta/\omega + t - 2\zeta/\omega = t$. The homogeneous solutions are $e^{-\zeta\omega t}(A\cos\wdm t +
B\sin\wdm t)$. The initial conditions give
\begin{align}
  \rho(0) &= -\frac{2\zeta}{\omega^3} + A = 0 &&\Rightarrow\quad A = \frac{2\zeta}{\omega^3},\\
  \dot\rho(0) &= \frac{1}{\omega^2} - \zeta\omega A + \wdm B = 0 &&\Rightarrow\quad
  B = \frac{\zeta\omega A - 1/\omega^2}{\wdm} = \frac{2\zeta^2-1}{\omega^2\wdm}.
\end{align}
Hence, for $t\ge0$,
\begin{equation}
  \rho(t) = \frac{t}{\omega^2} - \frac{2\zeta}{\omega^3}
  + e^{-\zeta\omega t}\left[\frac{2\zeta}{\omega^3}\cos\wdm t + \frac{2\zeta^2-1}{\omega^2\wdm}\sin\wdm t\right].
\end{equation}

\subsection*{B.4 $z$-transforms}

Write $E = e^{-\zeta\omega\Delta t}$, $c = \cos\wdm\Delta t$, $s = \sin\wdm\Delta t$ and $w = q^{-1}$.
The transform of a sequence $u_l$ ($l\ge0$) is $U(q) = \sum_{l\ge0}u_lq^{-l}$. Four transforms are
needed.
\begin{itemize}[nosep]
\item $u_l = 1$: the geometric series $\sum_lw^l = \frac{1}{1-w}$.
\item $u_l = l\Delta t$: differentiating the geometric series, $\sum_l lw^l = w\frac{\dd}{\dd w}\frac{1}{1-w} =
\frac{w}{(1-w)^2}$, so the transform is $\frac{\Delta t\,w}{(1-w)^2}$.
\item $u_l = E^l\cos(l\wdm\Delta t)$ and $u_l = E^l\sin(l\wdm\Delta t)$: they are the real and imaginary
parts of $(Ee^{\ii\wdm\Delta t})^l$, whose transform is $\frac{1}{1-Ee^{\ii\wdm\Delta t}w}$. Multiplying
numerator and denominator by $1-Ee^{-\ii\wdm\Delta t}w$, the denominator becomes
$(1-Ee^{\ii\wdm\Delta t}w)(1-Ee^{-\ii\wdm\Delta t}w) = 1-2Ecw+E^2w^2 =: \mathcal D(w)$ and the
numerator $1-Ee^{-\ii\wdm\Delta t}w = 1 - Ecw + \ii Esw$. Therefore the transforms are
$\frac{1-Ecw}{\mathcal D(w)}$ and $\frac{Esw}{\mathcal D(w)}$.
\end{itemize}
Applying them to $\rho_l$:
\begin{equation}
  R(q) = \frac{\Delta t\,w}{\omega^2(1-w)^2} - \frac{2\zeta}{\omega^3}\,\frac{1}{1-w} + \frac{\mathcal N(w)}{\mathcal D(w)},
  \qquad \mathcal N(w) = n_0 + n_1w,
\end{equation}
with
\begin{equation}
  n_0 = \frac{2\zeta}{\omega^3},\qquad n_1 = -\frac{2\zeta}{\omega^3}Ec + \frac{2\zeta^2-1}{\omega^2\wdm}Es .
\end{equation}

\subsection*{B.5 The transfer function}

By~\eqref{eq:B2}, the transform of $\gamma_l$ is $\Gamma_d(q) = \frac{1}{\Delta t}(q-2+q^{-1})R(q)$
(a shift by one sample to the left multiplies a transform by $q$, to the right by $q^{-1}$). Since
$q-2+q^{-1} = q(1-w)^2$:
\begin{equation}
  \Delta t\,\Gamma_d(q) = q(1-w)^2R(q) = \frac{\Delta t}{\omega^2} - \frac{2\zeta}{\omega^3}(q-1) + q(1-w)^2\frac{\mathcal N(w)}{\mathcal D(w)},
\end{equation}
where $qw = 1$ was used in the first two terms ($q(1-w)^2\cdot\frac{w}{(1-w)^2} = 1$ and
$q(1-w)^2\cdot\frac{1}{1-w} = q(1-w) = q-1$). Over the common denominator $\mathcal D$, the numerator is
\begin{equation}
  \frac{\Delta t}{\omega^2}\mathcal D - \frac{2\zeta}{\omega^3}(q-1)\mathcal D + (q-2+w)(n_0+n_1w).
\end{equation}
The two products containing $q$ expand as follows (using $q\mathcal D = q - 2Ec + E^2w$):
\begin{align}
  (q-1)\mathcal D &= q - (1+2Ec) + (E^2+2Ec)w - E^2w^2,\\
  (q-2+w)(n_0+n_1w) &= n_0q + (n_1-2n_0) + (n_0-2n_1)w + n_1w^2 .
\end{align}
The terms in $q$ are $-\frac{2\zeta}{\omega^3}q + n_0q = 0$: they cancel, and the filter is causal.
Collecting the powers of $w$:
\begin{align}
  w^0:&\quad \frac{\Delta t}{\omega^2} + \frac{2\zeta}{\omega^3}(1+2Ec) + n_1 - 2n_0
  = \frac{\Delta t}{\omega^2} + \frac{2\zeta}{\omega^3}(Ec-1) + \frac{2\zeta^2-1}{\omega^2\wdm}Es,\\
  w^1:&\quad -\frac{2Ec\,\Delta t}{\omega^2} - \frac{2\zeta}{\omega^3}(E^2+2Ec) + n_0 - 2n_1
  = -\frac{2Ec\,\Delta t}{\omega^2} + \frac{2\zeta}{\omega^3}(1-E^2) - 2\,\frac{2\zeta^2-1}{\omega^2\wdm}Es,\\
  w^2:&\quad \frac{E^2\Delta t}{\omega^2} + \frac{2\zeta}{\omega^3}E^2 + n_1
  = \frac{E^2\Delta t}{\omega^2} + \frac{2\zeta}{\omega^3}(E^2 - Ec) + \frac{2\zeta^2-1}{\omega^2\wdm}Es .
\end{align}
(In the $w^0$ line: $\frac{2\zeta}{\omega^3}(1+2Ec) - 2n_0 - \frac{2\zeta}{\omega^3}Ec =
\frac{2\zeta}{\omega^3}(1+2Ec-2-Ec) = \frac{2\zeta}{\omega^3}(Ec-1)$. In the $w^1$ line the terms
$\mp\frac{4\zeta}{\omega^3}Ec$ cancel.) Dividing by $\Delta t$ and writing every coefficient over
$\omega^3\Delta t$:
\begin{result}[Ramp-invariant coefficients]
\begin{align}
  b_1 &= \frac{\omega\Delta t + 2\zeta(Ec-1) + \frac{\omega}{\wdm}(2\zeta^2-1)Es}{\omega^3\Delta t},\\
  b_2 &= \frac{-2\omega\Delta t\,Ec + 2\zeta(1-E^2) - 2\frac{\omega}{\wdm}(2\zeta^2-1)Es}{\omega^3\Delta t},\\
  b_3 &= \frac{(2\zeta+\omega\Delta t)E^2 + E\bigl(\frac{\omega}{\wdm}(2\zeta^2-1)s - 2\zeta c\bigr)}{\omega^3\Delta t},\\
  a_1 &= -2Ec,\qquad a_2 = E^2,
\end{align}
\end{result}
and $\Gamma_d(q) = (b_1+b_2q^{-1}+b_3q^{-2})/(1+a_1q^{-1}+a_2q^{-2})$. By~\eqref{eq:B1} the sequence
$z_l$ is the input sequence filtered by $\Gamma_d$, which is the recursion~\eqref{eq:recursion}.

\emph{Check.} At zero frequency ($q=1$) the gain must be the static deflection $1/\omega^2$.
Numerically, for $\fn = 20$\,Hz, $Q=10$, $f_s = 500$\,Hz, $(b_1+b_2+b_3)/(1+a_1+a_2) = 6.33257\times10^{-5}$
and $1/\omega^2 = 6.33257\times10^{-5}$. Inserting $q = e^{\ii\theta}$ gives the frequency
response~\eqref{eq:Hd}; the program evaluates it in closed form for all nodes and frequencies at once
and checks it against a direct evaluation of the polynomials.

\section{Results quoted from the literature}\label{app:quoted}

\subsection*{C.1 Dirlik's distribution of rainflow ranges}\label{app:dirlik}

Dirlik fitted, on a large set of simulated Gaussian processes, a closed-form probability density of
the rainflow ranges $S$ of a stationary Gaussian process with the spectral moments $m_0,m_1,m_2,m_4$
(Section~\ref{sec:respvar}). With
\begin{equation}
  x_m = \frac{m_1}{m_0}\sqrt{\frac{m_2}{m_4}},\qquad \gamma = \frac{m_2}{\sqrt{m_0m_4}},
\end{equation}
\begin{gather}
  D_1 = \frac{2(x_m-\gamma^2)}{1+\gamma^2},\qquad
  R = \frac{\gamma - x_m - D_1^2}{1-\gamma-D_1+D_1^2},\qquad
  D_2 = \frac{1-\gamma-D_1+D_1^2}{1-R},\\
  D_3 = 1-D_1-D_2,\qquad
  Q_D = \frac{1.25(\gamma-D_3-D_2R)}{D_1},
\end{gather}
the density of the normalised range $Z = S/(2\sqrt{m_0})$ is
\begin{equation}
  p(Z) = \frac{D_1}{Q_D}e^{-Z/Q_D} + \frac{D_2\,Z}{R^2}e^{-Z^2/(2R^2)} + D_3\,Z\,e^{-Z^2/2},
\end{equation}
and the cycles occur at the rate of the peaks, $\nu_p = \sqrt{m_4/m_2}$. The damage uses the amplitude
$A = S/2 = \sqrt{m_0}\,Z$. Its moment of order $b$ follows term by term: for the exponential term the
substitution $u = Z/Q_D$ gives $\int_0^\infty Z^b\frac{D_1}{Q_D}e^{-Z/Q_D}\dd Z = D_1Q_D^b\Gamma(1+b)$;
the other two terms are Rayleigh densities with $\sigma = R$ and $\sigma=1$ (weighted by $D_2$ and $D_3$),
whose moments are given by~\eqref{eq:rayleighmoment}. Hence the expected damage per second
\begin{equation}
  d_{\mathrm{Dir}} = \frac{K^b}{C}\,\nu_p\,m_0^{b/2}\Bigl[D_1Q_D^b\,\Gamma(1+b) +
  2^{b/2}\,\Gamma\!\left(1+\tfrac b2\right)\bigl(D_2|R|^b + D_3\bigr)\Bigr].
\end{equation}
For a narrow-band process ($\gamma\to1$) the formula tends to the Rayleigh result $\nu_0^+(2m_0)^{b/2}
\Gamma(1+b/2)$ of~\eqref{eq:dgauss} (checked numerically in the tests within 2\,\% for a pure tone
spectrum). The program evaluates the formula with safeguards against divisions by zero for nearly
pure tones.

\subsection*{C.2 Central limit theorem}

A sum of many independent random variables with finite variances, none of which dominates the sum,
is approximately Gaussian. Section~\ref{sec:multisine} applies it to the lines of a random-phase
multisine.

\subsection*{C.3 Welch's PSD estimate}

Welch's estimate averages the periodograms of overlapping windowed segments. Its expectation is the
true PSD smoothed over the resolution of a segment, and its relative scatter decreases as the square
root of the number of independent segments. The program uses segments of length equal to the inverse of
the requested resolution, a Hamming window and 50\,\% overlap.


\end{document}
